\documentclass[conference]{IEEEtran}

\usepackage[T1]{fontenc}
\usepackage[utf8]{inputenc}
\usepackage[brazilian]{babel}
\usepackage{microtype}
\usepackage{booktabs}
\usepackage{array}
\usepackage{graphicx}
\usepackage{flafter}
\usepackage{float}
\usepackage{cite}
\usepackage{url}
\usepackage[hidelinks]{hyperref}
\usepackage{tikz}
\usetikzlibrary{arrows.meta,positioning,shapes.geometric,calc}

\newcommand{\system}{\textsc{Cognalys}}
\newcommand{\researchq}[1]{\textbf{QP#1}}
\newcommand{\pending}[1]{\textbf{[TODO: #1]}}

\title{Cognalys: Uma Biblioteca Python que Complementa o RVSec com Análise e Reparo Guiados por Evidências}

\author{\IEEEauthorblockN{João Pedro Antonow}
\IEEEauthorblockA{Departamento de Ciência da Computação (CIC)\\
Universidade de Brasília (UnB)\\
Brasília, Brasil\\
\texttt{jpantonow@gmail.com}\\[3pt]
Orientador: Prof.\ Rodrigo Bonifácio\\Coorientador: Prof.\ Guilherme Novaes Ramos}}

\begin{document}
\renewcommand{\IEEEkeywordsname}{Palavras-chave}

\maketitle

\begin{abstract}
O RVSec relata usos indevidos de APIs criptográficas que uma aplicação Android executou,
mas não qual alteração no código-fonte é segura: uma correção pode depender de dados armazenados, de
pares de protocolo, de chamadores e da política da aplicação. O \system, uma biblioteca Python e
ferramenta de linha de comando, trata cada achado do RVSec como uma afirmação a ser verificada antes
de corrigi-lo. Com evidências do repositório na revisão analisada, etapas baseadas em modelos
de linguagem de grande porte, sob verificações determinísticas, decidem se a violação
é real e se uma correção preserva o comportamento da aplicação. Uma
violação confirmada com tal correção, mesmo que seja apenas de regra, torna-se um patch aplicado em um
worktree isolado na ponta do branch alvo e submetido a build; os demais achados
terminam com seu motivo registrado, ou com uma issue quando seu impacto de segurança
foi confirmado. Nada é publicado sem aprovação humana explícita.
Relatamos um estudo de caso datado de uma versão anterior do pipeline e duas
avaliações pré-registradas com desvios declarados; todos os resultados são
descritivos. Em 18 de novembro de 2026, dos 49 pull requests submetidos para achados
do RVSec da campanha do grupo do RVSec (8), de relatos anteriores do RVSec (15)
e de nossas próprias execuções do RVSec (26), 21 haviam sido incorporados (um apenas adiciona um teste), 10 haviam sido fechados e 18
estavam abertos; seis dos dez fechamentos questionaram a relevância de segurança do
uso relatado. Quando o RVSec foi reexecutado sobre o APK original e o APK com patch, a
violação relatada não apareceu mais no método alterado pelo patch em 33 dos
38 pull requests que ele pôde verificar (25 pela regra registrada), e um patch corrigido manteve
o mesmo material de chave e apenas mudou a forma como o monitor o via. Em 20 unidades de análise, dois juízes LLM, sem saber qual braço produziu cada
entrega, compararam o \system{} com agentes de programação em dois blocos. Agentes
que receberam apenas o achado e o repositório propuseram menos pull requests do
que o \system{}, com menos julgados aceitáveis (7 a 18, contra 28). Agentes
instruídos a seguir o mesmo processo propuseram mais, com tantos ou mais julgados
aceitáveis na leitura completa (28 e 32), mas também mais julgados inaceitáveis
(10 e 14, contra 3). Um escore por execução que conta issues e desfechos de nenhuma
ação como aceitáveis, e assim recompensa a abstenção, colocou o primeiro bloco
perto do \system{} ou acima dele (0,90 a 0,97, contra 0,96) e o segundo abaixo
dele (0,85 e 0,80). Nem o build nem a reexecução mostram que um patch preserva o
comportamento da aplicação.
\end{abstract}

\begin{IEEEkeywords}
uso indevido de APIs criptográficas, verificação em tempo de execução, Android, modelos de linguagem
de grande porte, reparo de programas, agentes de programação, contribuições de código aberto, falsos
positivos
\end{IEEEkeywords}

\section{Introdução}
\label{sec:introduction}

Desenvolvedores têm dificuldade em usar corretamente as APIs criptográficas
de Java~\cite{nadi2016hoops}, e usos indevidos continuam comuns em aplicações
Android~\cite{egele2013crypto,costa2026rv4android}; muitas ferramentas os
detectam: analisadores estáticos como CogniCrypt e CryptoGuard, e ferramentas dinâmicas
como CryLogger e RVSec~\cite{krueger2017cognicrypt,krueger2021crysl,
rahaman2019cryptoguard,piccolboni2021crylogger,torres2023runtime}. A detecção foi
amplamente estudada; agir a partir de um relato, muito menos. Um relato diz que uma
regra foi violada, não se a violação importa naquela aplicação nem como
alterar o programa sem quebrá-lo. Substituir \texttt{SHA-1} por
\texttt{SHA-256} altera os bytes que o armazenamento, a sincronização ou um servidor
esperam, e fixar uma versão de TLS remove protocolos de que um par ainda pode precisar. Os
mantenedores que responderam aos nossos pull requests levantaram exatamente essas preocupações:
comportamento que consideram intencional, a compatibilidade de valores armazenados e
se a regra importa em seu contexto (Seção~\ref{sec:results}).

O RVSec~\cite{torres2023runtime,costa2026rv4android} monitora chamadas a APIs criptográficas enquanto um
APK instrumentado executa e relata cada violação no método que fez a
chamada. Seus achados mostram que o uso indevido é executado, o que um alerta estático não
mostra, mas referem-se a um APK compilado e não ao código-fonte: o método
relatado pode ser uma lambda ou corrotina Kotlin, um nome alterado pelo compilador ou
código de biblioteca empacotado no aplicativo. Transformar tal achado em uma alteração que um
mantenedor possa aceitar exige mais três passos: localizar o código na
revisão do repositório a partir da qual o APK foi gerado, decidir se a violação é
real e se importa ali, e produzir um patch que preserve o comportamento da aplicação
e passe no build.

Modelos de linguagem de grande porte tornam possível tentar esses passos: eles corrigem código,
às vezes melhor do que sistemas de reparo
anteriores~\cite{xia2023apr,zhang2023apimisuse}, e agentes de programação resolvem parte das issues reais
de repositórios, como mede o benchmark
SWE-bench~\cite{jimenez2024swebench,yang2024sweagent}. Eles também tornam os passos mais arriscados: código gerado pode conter
fraquezas de segurança~\cite{pearce2022asleep}, um patch pode satisfazer seu oráculo, aqui o
build, sem estar correto~\cite{smith2015overfitting}, e um modelo pode
errar tanto sobre se uma violação importa quanto sobre seu próprio patch. Uma
pergunta natural é se um agente de programação de propósito geral basta, dado o
achado e o repositório, ou dado também o mesmo processo escrito como
instruções (Seção~\ref{subsec:rq6}).

Este artigo apresenta o \system, uma biblioteca Python e ferramenta de linha de comando para aplicações
Android em Java e Kotlin que verifica cada achado de uso indevido de API antes de
agir sobre ele. Para cada achado, ele decide, em ordem, se a
violação relatada de fato ocorre no código e se existe uma correção que preserve o
comportamento da aplicação. Uma violação confirmada com tal correção torna-se um patch,
tenha ela impacto de segurança ou não. Caso contrário, o achado termina com seu
motivo registrado: um falso positivo, código que já usa a API conforme documentado
(de modo que qualquer correção mudaria aquilo em que ele confia ou o que protege), ou uma correção que exige
uma decisão fora do repositório; uma violação confirmada sem correção torna-se
uma issue somente quando seu impacto de segurança foi confirmado. Essas abstenções são
desfechos válidos, não falhas. A principal escolha de projeto é separar
responsabilidades: componentes determinísticos coletam evidências, etapas de LLM
adjudicam, planejam, redigem, revisam e descrevem o patch sob esquemas
versionados, e validadores determinísticos rejeitam afirmações sem suporte, planos
incompletos, padrões de falha conhecidos e textos que descrevem edições ausentes do
diff. Patches aceitos são portados para a ponta do branch alvo e passam por build
antes que alguém seja solicitado a publicá-los, e as respostas de mantenedores a
pull requests anteriores são lidas por análises posteriores. Achados de outros detectores nos
formatos SARIF ou Semgrep JSON entram no mesmo fluxo de trabalho, mas não são avaliados
aqui.

Este trabalho traz cinco contribuições:
\begin{itemize}
  \item uma biblioteca Python e CLI que complementam o RVSec ao conectar cada
  achado de tempo de execução em um APK à revisão do código-fonte de onde ele veio, a uma decisão
  sobre se ele é real e a um patch ou issue passível de revisão;
  \item um projeto em etapas com contratos separados para evidências, adjudicação,
  edições autorizadas, validação isolada e revisão humana, no qual a abstenção
  e toda decisão intermediária permanecem inspecionáveis;
  \item um estudo dos 49 pull requests (46 sobre especificações criptográficas)
  e 10 issues submetidos a mantenedores até 18 de novembro de 2026 para achados
  do RVSec da campanha do grupo do RVSec (8 pull requests), de relatos anteriores
  do RVSec (15) e de nossas próprias execuções do RVSec (26), com 21 pull requests
  incorporados, 10 fechados e 18 abertos, incluindo os motivos que os mantenedores
  deram para suas decisões;
  \item uma validação dinâmica, pré-registrada antes do build de qualquer caso novo, na
  qual o RVSec reexecuta o APK original e o APK com patch de 45 pull requests
  publicados (33 que atenderam aos critérios de inclusão e 12 dos 13 acrescentados
  após seus resultados): a violação relatada não apareceu mais em 33 dos 38 que ele
  pôde verificar (25 pela regra registrada); e
  \item uma comparação pré-registrada em 20 unidades de análise (24 achados) com
  agentes de programação (Codex~\cite{openai2026codex}
  com dois modelos e Claude Code~\cite{anthropic2026claudecode}) em dois blocos, que receberam o achado e o
  repositório, ou também o mesmo processo com o porte do patch para a ponta do
  branch alvo, com as entregas pontuadas por dois juízes LLM: os agentes do primeiro
  bloco propuseram menos pull requests do que o \system{}, com menos julgados aceitáveis;
  os agentes do segundo bloco tiveram tantos ou mais julgados aceitáveis na leitura
  completa, mas mais julgados inaceitáveis (10 e 14, contra 3). Com o mesmo modelo e as mesmas
  instruções do \system{} em um dos braços, a diferença é compatível com a
  execução do processo em etapas separadas, com a decisão de agir tomada à parte
  da edição, e não com as instruções.
\end{itemize}

A Seção~\ref{sec:background} revisa os trabalhos relacionados,
a Seção~\ref{sec:pipeline} descreve o pipeline, a Seção~\ref{sec:settings}
as questões de pesquisa, um estudo de caso do pipeline como estava em outubro
de 2026 e as duas avaliações pré-registradas, a Seção~\ref{sec:results} os
resultados, e a Seção~\ref{sec:discussion} as implicações e as ameaças à
validade; a Seção~\ref{sec:conclusion} conclui.
\section{Fundamentação e Trabalhos Relacionados}
\label{sec:background}

\subsection{Usos Indevidos de APIs Criptográficas}

Uma especificação de uso de API descreve restrições que vão além da verificação
de tipos, como chamar \texttt{load} antes de ler um repositório de chaves, usar
um digest adequado, inicializar uma cifra de forma consistente com seu modo ou
impedir a negociação de protocolos obsoletos; algumas dependem de fluxo de
dados, do comportamento do provedor ou do estado em tempo de execução.
Analisadores estáticos como CogniCrypt
e CryptoGuard verificam código-fonte ou bytecode contra essas regras, e o
CryLogger observa valores em tempo de
execução~\cite{krueger2017cognicrypt,rahaman2019cryptoguard,piccolboni2021crylogger}.
CrySL, a linguagem de regras por trás do CogniCrypt, combina sequências de
eventos permitidas, restrições sobre valores e predicados que um objeto repassa
a outro~\cite{krueger2021crysl}.

O RVSec aplica verificação em tempo de execução ao mesmo problema; suas
especificações criptográficas são monitores JavaMOP~\cite{jin2012javamop}
traduzidos manualmente de regras CrySL, e técnicas estáticas e de tempo de
execução revelam conjuntos diferentes de usos indevidos e são
complementares~\cite{torres2023runtime}. O RVSec é a ferramenta do grupo,
publicada como RV4Android~\cite{costa2026rv4android}: ele instrumenta com esses
monitores o bytecode de um APK e exercita o aplicativo instrumentado em um
emulador com ferramentas de geração de testes como Monkey~\cite{androidmonkey},
APE~\cite{gu2019ape} e DroidBot~\cite{li2017droidbot}. Nos relatórios que
analisamos, cada violação traz um código que nomeia a restrição violada e,
quando disponíveis, os valores observados e esperados. O conjunto criptográfico
na versão que usamos (jca\_android) tem 47 especificações que seguem as regras
CrySL. Executado com onze configurações de exploração sobre 163 aplicações
do F-Droid~\cite{fdroid}, o RVSec relatou 453 usos indevidos distintos em 89 aplicações, 82\% deles
em código de bibliotecas embutidas, enquanto os métodos que chamam diretamente uma
API monitorada raramente foram exercitados~\cite{costa2026effectiveness}. Aquele
estudo mede o que o RVSec encontra; este trabalho parte de seus achados e de nossas
próprias execuções do RVSec, e pergunta que mudança, se alguma, um mantenedor
deveria receber.

Cada fonte de relato traz sua própria incerteza. Um monitor vê apenas os
caminhos alcançados pela exploração, de modo que um uso indevido por trás de uma
tela de login ou de configuração não é relatado, e ferramentas de exploração
diferentes alcançam códigos diferentes~\cite{costa2026effectiveness}. Uma regra pode ser mais rígida que o
comportamento documentado da plataforma: o nome de um algoritmo do Android
Keystore pode estar ausente de uma regra escrita para provedores de desktop~\cite{costa2026effectiveness}.
Alguns relatos dizem que um valor \emph{não foi observado} como proveniente de uma fonte
aprovada, como uma chave gerada fora da execução monitorada: um limite do
rastreamento, não um uso indevido. Na campanha do grupo, objetos inicializados com
padrões da plataforma que o monitor nunca vê respondem por 131 dos 453 usos
indevidos~\cite{costa2026effectiveness}. Detectores estáticos de uso indevido de API também apresentaram
baixa precisão e revocação~\cite{amann2019evaluation}. Um achado é, portanto, evidência, e não
prova de que uma determinada edição está correta.
O \system{} transforma três dessas incertezas em perguntas que responde antes
de qualquer edição: mapeia o método relatado de volta ao código-fonte, confronta
a regra com o comportamento documentado da plataforma e permite que a
adjudicação conclua que um limite de rastreamento não é um uso indevido
(Seção~\ref{sec:pipeline}). A cobertura continua sendo
o limite do detector.

\subsection{Reparo Automatizado e Agentes de Programação}
\label{subsec:apr}

O reparo automatizado de programas (APR) produz patches candidatos para um
defeito relatado e mantém aqueles que um oráculo, geralmente os testes do
projeto, aceita~\cite{legoues2019apr}; um patch pode satisfazer esse oráculo sem
corrigir o defeito, problema conhecido como sobreajuste de
patch~\cite{smith2015overfitting}. O \system{} não dispõe de teste que codifique
o uso indevido, e nenhuma das verificações que aceitam seus patches observa se a
violação desapareceu (Seção~\ref{sec:pipeline}); a avaliação, em vez disso,
reexecuta o detector (Seção~\ref{subsec:dynval}). Onde o reparo foi implantado em escala, uma pessoa
permanece no circuito: no Facebook, o SapFix e o Getafix propunham correções que
engenheiros revisavam antes da integração~\cite{marginean2019sapfix,bader2019getafix}.
O \system{} mantém esse modelo: cada decisão registra sua justificativa, e o
autor aprova cada submissão.

Ferramentas de reparo de usos indevidos de criptografia corrigem o que um
detector relata. O CDRep detecta sete tipos de uso indevido em bytecode Android
e reescreve cada um com um padrão de correção~\cite{ma2016cdrep}; o Seader
infere modelos de detecção e de reparo para as APIs criptográficas de Java a
partir de pares de código inseguro e seguro~\cite{zhang2022seader}. O reparo
baseado em LLM~\cite{xia2023apr,zhang2023apimisuse} também parte de um defeito
conhecido; no reparo zero-shot de falhas de segurança, modelos de linguagem
corrigiram vulnerabilidades sintéticas e construídas à mão, mas tiveram
dificuldade em produzir correções funcionalmente corretas para vulnerabilidades
reais~\cite{pearce2023zeroshot}. Todos tomam o defeito relatado como dado,
premissa que nossos mantenedores questionaram (\researchq{3}, \researchq{4}); o
\system{}, por isso, decide se um achado é real e se tem uma correção que
preserve o comportamento antes de gerar qualquer edição.

Agentes de programação executam um modelo em um laço com ferramentas, lendo e
editando arquivos e executando comandos até que o modelo decida que a tarefa
está concluída~\cite{yang2024sweagent}; o Codex CLI~\cite{openai2026codex} e o
Claude Code~\cite{anthropic2026claudecode} executam esse laço na máquina de um
desenvolvedor. O SWE-bench aceita um patch para uma issue real do GitHub quando
passam os testes extraídos do pull request que a resolveu~\cite{jimenez2024swebench},
um oráculo sujeito ao mesmo sobreajuste. A estrutura do trabalho importa: o
SWE-agent mostrou que a interface pela qual um agente visualiza e edita um
repositório altera quantas issues ele resolve~\cite{yang2024sweagent}, e o
Agentless, uma sequência fixa de localização, reparo e validação de patch em que
nenhum modelo escolhe a próxima ação, superou os agentes de código aberto então
disponíveis no SWE-bench Lite a um custo menor~\cite{xia2025agentless}. O
Agentless compara estruturas diante de um oráculo de testes; a tarefa do
\system{} inclui também decidir se deve agir. Um agente segue um processo apenas
na medida em que segue suas instruções, ao passo que o \system{} executa cada
etapa em código e verifica sua saída antes que a próxima comece. A comparação da
Seção~\ref{subsec:protocol} fornece aos agentes de programação o processo do
\system{} como instruções, incluindo o porte para a ponta do branch, e
assim mede a diferença entre um processo que é instruído e outro que é imposto.

\subsection{Contribuições Geradas por Ferramentas para Código Aberto}
\label{subsec:bots}

Em um estudo de projetos populares do GitHub, cerca de um quarto usava bots, e
os mantenedores os consideraram úteis, porém limitados~\cite{wessel2018bots}.
Pull requests automatizados que atualizam dependências desatualizadas levaram os
projetos a atualizá-las com mais frequência, mas muitos ficaram sem merge, e os
desenvolvedores citaram o risco de mudanças incompatíveis e o esforço de
compreender a
mudança~\cite{mirhosseini2017autopr}. Pull requests abertos por bots são
mergeados com menos frequência e esperam mais por uma resposta do que os abertos
por pessoas~\cite{wyrich2021botprs}. Uma atualização de dependência pode ser
verificada contra uma release; uma correção de uso indevido de API altera o
próprio código do projeto, e se ela é desejada depende do que esse código
protege. Todo pull request e toda issue deste estudo foram aprovados pelo autor
antes de serem abertos; a
Seção~\ref{subsec:form} mostra como os mantenedores reagiram à forma das
primeiras submissões e como o texto mudou.

\subsection{LLMs como Juízes}
\label{subsec:judges-bg}

Em perguntas de assistentes de conversação, juízes LLM fortes concordaram com
avaliadores humanos com frequência semelhante à concordância entre os próprios
humanos, mas foram influenciados pela posição e pelo tamanho das respostas e
podem favorecer suas próprias saídas~\cite{zheng2023judge}. A ordem em que os candidatos são
apresentados pode mudar um veredito~\cite{wang2024fair}, e avaliadores LLM
reconhecem e favorecem suas próprias
gerações~\cite{panickssery2024selfpref}. Em artefatos de engenharia de software,
as avaliações de LLMs se aproximam da concordância humana em algumas tarefas de
anotação, mas não em outras, e a concordância entre modelos sinaliza quando
eles são
confiáveis~\cite{ahmed2025annotation}. A comparação da
Seção~\ref{subsec:protocol} usa dois juízes de provedores diferentes e trata
esses vieses: os cartões são pontuados em ordem aleatória; uma das leituras
mostra apenas o patch, sem seu título e descrição; a concordância entre
os juízes é relatada ao lado da pontuação final; e uma sondagem pede a um juiz que
identifique o braço por trás de um cartão; a Seção~\ref{subsec:threats}
discute o risco que permanece.

\subsection{Posicionamento Deste Trabalho}

Detectores, estáticos (CogniCrypt, CryptoGuard) ou dinâmicos (CryLogger,
RVSec), terminam com um relato; o \system{} o trata como uma alegação a ser
verificada contra a revisão do código-fonte, sem presumir impacto de segurança,
e decide se e como alterar o código. O reparo baseado em LLM, incluindo o
trabalho mais próximo, um preprint de Zhang et al.\ que avalia modelos de
linguagem pré-treinados no reparo de usos indevidos de
API~\cite{zhang2023apimisuse}, pontua patches candidatos para defeitos
conhecidos em benchmarks; o \system{} pode se abster, autoriza edições por
arquivo e símbolo, faz o build de cada patch no projeto real e, neste estudo,
enviou seus patches a mantenedores reais e os teve verificados novamente pelo
RVSec. O SWE-bench e o SWE-agent compartilham o cenário em nível de
repositório, mas aceitam um patch quando passam os testes da issue, ao passo
que o \system{} parte de evidências do analisador e verifica um patch com um
build, uma regra estática quando existe e uma revisão do seu comportamento por
modelo. O Codex CLI e o Claude Code fazem edições em um laço de agente; são as
linhas de base da nossa comparação, que recebem o processo do \system{} como
instruções. Pull requests de bots trazem mudanças que um oráculo existente,
como uma release, pode verificar; as mudanças do \system{} alteram código
relevante para a segurança, e cada uma é aprovada por uma pessoa antes de ser
enviada.

\section{O Pipeline \system{}}
\label{sec:pipeline}
\label{subsec:library}

O \system{} é uma biblioteca Python e ferramenta de linha de comando que
transforma cada achado reportado por um detector em um desfecho revisado: uma
proposta de correção, uma issue para os mantenedores ou uma decisão documentada
de não propor nada (Figura~\ref{fig:pipeline}). Cada etapa grava seus
resultados no diretório da execução, em vez de passar uma conversa opaca à
etapa seguinte, de modo que toda decisão pode ser rastreada até as evidências
que utilizou.

O \system{} é construído em torno de sete propriedades, cada uma com um limite
declarado. \emph{Determinismo onde o código pode decidir:} a coleta de
evidências, as verificações de validação e o build são determinísticos; os
estágios de modelo que adjudicam, planejam, redigem e revisam não são, de modo
que o mesmo achado pode terminar de formas diferentes entre execuções
(Seção~\ref{subsec:threats}). \emph{Rastreabilidade:} cada estágio grava no
diretório da execução suas entradas, saídas e os identificadores de evidência
que cita, e uma afirmação que cita um identificador desconhecido é rejeitada.
\emph{Auditabilidade:} cada chamada de modelo é registrada com sua requisição
e resposta; prompts, schemas, políticas, especificações, construtores de
contexto e resolvedores têm versão registrada, cada requisição e cada pacote de
evidências têm hash, e cada execução registra o commit da biblioteca e um hash
de seu código-fonte. \emph{Abstenção como desfecho:} um mapeamento ambíguo
fica sem resolução em vez de ser adivinhado, e terminar sem patch, com o
motivo registrado, é um resultado válido. \emph{Autoridade apenas da
evidência:} pesquisa externa, respostas de mantenedores e saídas de modelo
informam decisões, mas só evidências do repositório e política explícita podem
autorizar uma edição, de modo que conteúdo da web não pode instruir o
pipeline. \emph{Isolamento:} os patches são aplicados e compilados em um
worktree local separado, nada chega ao projeto antes da publicação, e nenhuma
credencial entra em um prompt, embora a análise possa contatar hosts e
baixar arquivos que o código revisado nomeia, para verificar suas cadeias de
certificados ou digests esperados; os builds de terceiros ainda não rodam em
contêiner (Seção~\ref{subsec:ethics}). \emph{Autonomia limitada:} uma execução
segue sem supervisão de um achado até um patch compilado, mas nada é publicado
até que uma pessoa aprove aquele achado específico.

\begin{figure*}[t]
\centering
\resizebox{\textwidth}{!}{%
\begin{tikzpicture}[
  box/.style={draw,rounded corners,align=center,minimum height=11mm,
              text width=2.5cm,font=\footnotesize},
  opt/.style={box,dashed},
  inp/.style={draw,dashed,rounded corners,align=center,minimum height=8mm,
              text width=2.6cm,font=\footnotesize,fill=black!4},
  gate/.style={draw,diamond,aspect=2,align=center,font=\scriptsize,
               inner sep=0pt,text width=1.3cm},
  stop/.style={draw,rounded corners,align=left,font=\footnotesize,
               text width=3.1cm,fill=black!8},
  flow/.style={-{Latex[length=2mm]},thick},
  fail/.style={-{Latex[length=2mm]},thick,densely dotted},
  feed/.style={-{Latex[length=2mm]},dashed},
  lbl/.style={font=\scriptsize,align=center},
  lane/.style={font=\footnotesize\bfseries,rotate=90,anchor=south}
]
\def\cx#1{#1*2.95}
% ---------------- lane 1: analysis
\node[opt] (rvrun)  at (\cx{0},0) {Opcional: APK compilado sem minificação; execução do RVSec (Monkey, APE)};
\node[box] (reports) at (\cx{1},0) {Relatórios do RVSec: rótulos interpretados, ecos agrupados, código do app ou de biblioteca};
\node[box] (map)    at (\cx{2},0) {Clone; nome de bytecode para a declaração no código-fonte; índices do código};
\node[box] (units)  at (\cx{3},0) {Unidades de análise agrupadas; pacote de evidências determinístico};
\node[box] (claims) at (\cx{4},0) {Extração de afirmações e contrato da função};
\node[gate] (adj)   at (\cx{5},0) {Adjudicação};
\draw[flow] (rvrun)--(reports);
\draw[flow] (reports)--(map);
\draw[flow] (map)--(units);
\draw[flow] (units)--(claims);
\draw[flow] (claims)--(adj);
% inputs
\node[inp] (research) at (\cx{3},1.75) {Pesquisa externa (fontes oficiais; parte compartilhada vinda do corpus)};
\node[inp] (spec)     at (\cx{4.5},1.75) {Especificação da regra, auditada contra a regra};
\node[inp] (memory)   at (\cx{6},1.75) {Respostas de mantenedores e patches anteriores};
\draw[feed] (research)--(units);
\draw[feed] (spec.south)--(claims.north);
\draw[feed] (spec.south)--(adj.north);
\draw[feed] (memory.south west)--(adj.north east);
% ---------------- outcome without patch
\node[stop,minimum height=26mm] (stop) at (\cx{6},-2.7) {\textbf{Termina sem patch, com o motivo registrado:} falso positivo, evidência insuficiente, sem impacto de segurança, sem correção segura, contrato externo, decisão do mantenedor, já proposto ou corrigido upstream};
\draw[fail] (adj.east) -| node[above,lbl,pos=0.25]{sem correção} (stop.north);
% ---------------- lane 2: repair
\node[box] (plan)   at (\cx{0},-3.2) {Plano do patch: edições, invariantes, mudanças proibidas};
\node[box] (drafts) at (\cx{1},-3.2) {Rascunhos por arquivo; verificações semânticas};
\node[gate] (review) at (\cx{2},-3.2) {Revisão comporta\-mental};
\node[box] (tip)    at (\cx{3},-3.2) {Worktree na ponta do branch: aplicação, merge de três vias ou porte pelo modelo};
\node[box] (build)  at (\cx{4},-3.2) {Build Gradle; reparo de compilação, revisado novamente};
\draw[flow] (adj.south) -- ($(adj.south)+(0,-0.75)$) node[right,lbl,pos=0.5]{correção} -| (plan.north);
\draw[flow] (plan)--(drafts);
\draw[flow] (drafts)--(review);
\draw[flow] (review)--node[above=7pt,xshift=-7pt,lbl]{limpo}(tip);
\draw[flow] (tip)--(build);
\draw[flow] (review.south west) to[out=-120,in=-60] node[below=-1pt,lbl]{regressão: novo rascunho\\ou novo plano (uma vez)} (drafts.south);
% failure exits (dotted) to the outcome box
\coordinate (bus) at ($(stop.west)+(0,0.85)$);
\draw[fail] (review.north) |- (bus);
\draw[fail] (tip.north) -- (tip.north |- bus);
\draw[fail] (build.north) -- (build.north |- bus);
\node[lbl,anchor=south west] at ($(review.north |- bus)+(0.05,0)$) {regressão persiste};
\node[lbl,anchor=south west] at ($(tip.north |- bus)+(0.05,0)$) {ponta já corrigida};
\node[lbl,anchor=south west] at ($(build.north |- bus)+(0.05,0)$) {build ainda falha};
% memory also informs planning and review
\draw[feed] (memory.east) -- ++(0.7,0) |- ($(plan.south)+(0,-1.12)$) -- (plan.south);
\draw[feed] ($(review.south)+(0.9,-1.12)$) -- ($(review.south)+(0.9,-0.3)$) -- (review.south east);
% ---------------- lane 3: proposal
\node[box] (text)    at (\cx{3},-6.1) {Texto do PR no template do repositório, ou texto da issue; conferido com o diff};
\node[box] (checks)  at (\cx{4},-6.1) {Verificações no repositório ao vivo: duplicata, mudança upstream, política};
\node[gate] (approve) at (\cx{5},-6.1) {Aprovação humana};
\node[box,text width=2.1cm] (publish) at (\cx{6}+0.35,-6.1) {Publicação remota explícita};
\draw[flow] (build.south) -- ($(build.south)+(0,-1.3)$) -| (text.north);
\draw[flow] (text)--(checks);
\draw[flow] (checks)--(approve);
\draw[flow] (approve)--node[above=7pt,lbl]{aprovado}(publish);
\draw[flow] (stop.south) -- ++(0,-0.35) -- ++(1.75,0) |- ($(text.south)+(0,-0.55)$) node[below,lbl,pos=0.75]{issue, apenas com impacto de segurança confirmado} -- (text.south);
% legend
\node[draw,rounded corners,align=left,font=\scriptsize,anchor=west,text width=6.3cm] at (\cx{0}-1.25,-6.1)
 {\tikz\draw[flow](0,0)--(0.6,0); fluxo principal\quad \tikz\draw[fail](0,0)--(0.6,0); termina sem patch\\[2pt]
  \tikz\draw[feed](0,0)--(0.6,0); informa uma decisão, não autoriza uma edição\\[2pt]
  caixa tracejada sombreada: entrada externa ao repositório};
% lane labels
\node[lane] at (\cx{0}-1.55,0) {Análise};
\node[lane] at (\cx{0}-1.55,-3.2) {Reparo};
\node[lane] at (\cx{0}-1.55,-6.1) {Proposta};
\end{tikzpicture}
}
\caption{Fluxo de trabalho exposto pela biblioteca Python e pela CLI. Setas
cheias são o fluxo principal; setas pontilhadas encerram uma unidade sem patch,
com o motivo registrado, e só um impacto de segurança confirmado transforma esse
fim em uma issue. Entradas tracejadas informam decisões, mas não podem autorizar
uma edição: só evidências do repositório e política explícita podem.}
\label{fig:pipeline}
\end{figure*}

\subsection{De um Relatório do RVSec ao Código-Fonte}

Uma tabela de dados no \system{} define o que cada rótulo do RVSec afirma: um
valor que a verificação não admite, chamadas fora da ordem da regra, uma origem
que o monitor não observou, uma origem observada que ele não admite, ou um eco
de um relatório sobre o produtor do objeto. Um eco é agrupado aos achados do
método que ele repete, a menos que sua origem esteja em outro lugar. Um achado
é marcado como código da aplicação ou de biblioteca a partir dos pacotes do
próprio app (os namespaces de todos os módulos do repositório); relatórios
sobre nomes de classe ofuscados são deixados de lado. Quando o \system{} inicia
a execução do RVSec, o APK é construído sem minificação onde a receita do
F-Droid permite, de modo que as localizações em tempo de execução correspondam
ao código-fonte. A observação em tempo de execução permanece nas evidências, e
seu rótulo delimita o que o relatório estabelece: uma origem não observada
indica o que o monitor viu, não o que o objeto é---um IV gerado pelo provedor
da cifra é aleatório, embora o monitor não consiga rastreá-lo.

O fluxo de trabalho clona o repositório, registra o branch base documentado, lê
seus arquivos de orientação e mapeia o método de bytecode para uma declaração
no código-fonte. O mangling de value classes do Kotlin e os sufixos do
compilador são removidos; o corpo de uma lambda ou corrotina que não nomeia
nenhuma função do código-fonte é mapeado para a única função de sua classe
externa que usa a API governada ou, na falta dela, para a única função que ele
chama no mesmo módulo Gradle e que a usa; um método ausente do arquivo
resolvido é procurado nos demais arquivos do pacote; e, em Kotlin
Multiplatform, o source set Android é preferido ao código desktop ou iOS que o
APK nunca contém. Uma correspondência ambígua permanece não resolvida em vez de
adivinhada. Índices conservadores de Java e Kotlin baseados em Tree-sitter~\cite{treesitter}
fornecem o método exato, o contexto da classe, chamadores, sobrescritas,
referências, consumidores de valores, testes, dependências, SDK mínimo e
configuração do Gradle; relações que exigem um classpath JVM completo
permanecem explicitamente ambíguas. O código-fonte alvo nunca é truncado; um
orçamento limita o contexto auxiliar e registra cada exclusão.

Achados em um mesmo trecho de código formam uma unidade com um único patch:
mesma classe e arquivo, e um método chama o outro, um método viola duas regras,
ou uma regra e uma mensagem se repetem ao longo do ciclo de vida de um objeto.
Achados de uma especificação em classes diferentes de um módulo também são
agrupados quando um método armazena um valor que o outro lê de volta, como um
IV gravado na cifragem e lido na decifragem. Rascunhos separados para esses
achados reescreviam o mesmo código de maneiras incompatíveis. Uma unidade
termina antes de qualquer chamada ao modelo quando cada achado já tem um pull
request nosso que está aberto ou foi mergeado após a revisão analisada, ou
quando o arquivo reportado não está no branch alvo.

Uma correção pode precisar ir além do método reportado. Para cada unidade, o
índice lista os arquivos que uma mudança pode afetar: chamadores em outros
arquivos, sobrescritas, consumidores de um símbolo compartilhado, testes
relacionados, configuração e, quando o valor protegido é armazenado, todo leitor
e escritor da mesma chave armazenada e todo arquivo por onde o valor passa. Uma
correção em vários arquivos é aplicada sem leitura humana
(\texttt{MULTI\_FILE\_FIX\_SUPPORTED}) somente quando esse escopo está completo,
a adjudicação revisou cada arquivo listado, todo arquivo que ela marca para
edição está no plano e o patch edita pelo menos dois arquivos. Fora disso, a
autoridade para aplicar só diminui: quando o índice foi truncado, um nome genérico ficou sem
resolução, um arquivo listado não foi revisado ou o plano edita além do método
reportado, a decisão passa a
\texttt{MULTI\_FILE\_FIX\_}\allowbreak\texttt{PROPOSED\_FOR\_REVIEW}, e os
arquivos não revisados são listados para a pessoa que a aprova; uma decisão que
não é apenas de revisão e não tem um arquivo exigido não gera patch. Cada
arquivo é redigido separadamente, diante da interface declarada dos demais, e
uma verificação que falha regenera apenas os arquivos que nomeia. A migração do
verificador armazenado do BitBanana mostra o risco: um candidato inicial
compilava sem atualizar leitores da credencial (\researchq{2}).

\subsection{Especificações e Pesquisa Externa}

Cada família de regras tem uma especificação JSON com os tipos governados, os
eventos, a semântica do detector e de tempo de execução, as formas conhecidas
de falso positivo e os requisitos de evidência, separando o que uma mensagem do
detector afirma do que estabeleceria um dano em tempo de execução (um aviso
literal de \texttt{TLS} não revela o protocolo negociado). As especificações
são escritas por um modelo a partir da regra e da documentação oficial e
auditadas deterministicamente contra a regra citada, de modo que um evento
omitido ou um valor acrescentado fica visível antes de influenciar uma análise.

Quatro especificações também trazem políticas de correção curadas; junto com
um arquivo do operador, que pode definir uma política por achado ou por família
de achados, elas são a política explícita que pode autorizar uma edição. Quando
uma política tem variantes, evidências determinísticas tiradas do código
revisado, e dos hosts e arquivos que ele nomeia, escolhem entre elas, e a
política resultante dá à adjudicação e ao planejamento o escopo, a estratégia,
as mudanças proibidas e os testes. Um plano completo em vários arquivos que
aguarda política só é aplicado sem revisão adicional quando a política o
autoriza e os dados protegidos se mostram restritos à memória; caso contrário,
é apenas de revisão.

A pesquisa externa opcional é uma chamada separada ao modelo, restrita a fontes
oficiais e primárias; o conteúdo da web ou do repositório é tratado como dado
não confiável, que não pode instruir o pipeline, escolher arquivos nem
autorizar publicação. Sua parte compartilhada, sobre a API e a regra, é mantida
em um corpus de pesquisa e reutilizada por até 90 dias apenas por uma
requisição idêntica byte a byte (mesmos prompts, esquema, modelo, esforço,
limite de saída e ferramenta de busca); a parte do caso é sempre pesquisada. O
estudo de caso e os pull requests publicados são anteriores ao corpus; as
execuções da comparação registram o estado a partir do qual começaram
(Seção~\ref{subsec:frozen}).

\subsection{Pacote de Evidências e Contratos com o LLM}

O Python constrói um pacote de evidências por unidade; cada item tem um
identificador, procedência, nível de confiança, estado de completude e hash.
Observações em tempo de execução e o código-fonte exato têm precedência sobre
relações semânticas, documentação oficial, pesquisa primária, fontes
secundárias e inferência do modelo. A ausência é
\texttt{not\_found\_in\_examined\_scope}, \texttt{confirmed\_absent} ou
\texttt{unknown\_due\_to\_}\allowbreak\texttt{incomplete\_scope}, valores nunca
usados de forma intercambiável. Seguem seis contratos com o modelo:
(1)~a \emph{extração de evidências} classifica afirmações como confirmadas,
inferidas, hipóteses ou desconhecidas, citando IDs de evidência, e declara o
contrato da função do código---seu propósito, sua contraparte (um serviço de
terceiros, um par operado pelo usuário, ou o próprio app ou dispositivo) e se
essa contraparte depende do comportamento reportado, citando literalmente o
código do repositório fora do método reportado; (2)~a \emph{adjudicação} decide
causa, violação, impacto de segurança em tempo de execução, compatibilidade,
escopo, bloqueadores de geração e de publicação, e a disposição;
(3)~o \emph{planejamento} autoriza cada arquivo, símbolo, tipo de edição,
invariante e teste, sem código; (4)~a \emph{geração} retorna substituições
exatas por arquivo; (5)~a \emph{revisão de comportamento} lê o patch finalizado
em relação ao arquivo inteiro e às invariantes e mudanças proibidas do plano, e
responde a cada preocupação da adjudicação; (6)~o \emph{texto da proposta}
descreve o diff final. Quando a extração aponta evidências ausentes ou a
adjudicação se abstém por falta delas, o Python as obtém e ambas são executadas
mais uma vez, mantendo a nova resposta somente se ela for validada. Etapas mais
estreitas, verificadas por esquema, tratam da pesquisa, do reparo de compilação
e do porte para a ponta do branch.

A adjudicação escolhe entre quinze decisões em cinco grupos: correções
aplicadas e construídas sem revisão adicional (correção em um ou em vários
arquivos, migração interna de digest, correção ou substituição de dependência);
correções que uma pessoa deve ler antes da aplicação (correção em vários
arquivos, migração de backend ou de vários repositórios); decisões que aguardam
uma política explícita; violações confirmadas que não exigem mudança aqui (uso
de digest sem fins de segurança, ausência de impacto de segurança, ou um
contrato externo de digest que nenhuma mudança no cliente consegue atender); e
\texttt{NO\_SAFE\_FIX}. Um patch de uma decisão apenas de revisão
ainda é compilado e revisado; ele só se torna publicável após um build aprovado,
um commit e uma revisão de comportamento atual sem regressão nem bloqueio e, como
toda proposta, apenas com a aprovação do autor.

Uma regra determinística também pode elevar uma resposta \texttt{NO\_SAFE\_FIX},
ou uma resposta em vários arquivos que aguarda política, à decisão apenas de
revisão \texttt{MULTI\_FILE\_FIX\_}\allowbreak\texttt{PROPOSED\_FOR\_REVIEW},
nunca a uma decisão aplicada sem revisão, quando a violação não foi julgada
falsa e cada bloqueador de geração é uma questão que quem lê o patch decide (a
autorização de um mantenedor, o comportamento no dispositivo, a compatibilidade
de dados armazenados ou de um consumidor) ou uma evidência que faltou ao
pacote. Uma causa que a leitura do código deixou em aberto passa então a se
apoiar na observação em tempo de execução ou, para um eco, no relatório de seu
produtor, e uma suposição de revisão registra isso; os bloqueadores de geração
restantes tornam-se bloqueadores de publicação, e o resultado deve passar pela
mesma validação. Uma correção confirmada recusada apenas por ter pedido geração
enquanto uma questão desse tipo seguia aberta também é mantida para revisão, em
vez de encerrada como \texttt{NO\_SAFE\_FIX}.

\subsection{Salvaguardas Determinísticas e Confiança Residual}

Verificações determinísticas protegem cada transição, porque uma saída
estruturada não torna uma conclusão verdadeira. Uma afirmação confirmada que
cita um ID de evidência desconhecido é rejeitada; uma adjudicação sustentada
exige uma causa e uma violação do detector confirmadas (em uma elevada à
decisão apenas de revisão, a causa pode se apoiar na observação em tempo de
execução ou no relatório de um produtor); um plano deve cobrir
todos os arquivos que o escopo exige, e cada arquivo gerado deve corresponder a
uma entrada do plano. Cada substituição deve ocorrer exatamente uma vez no
código-fonte base; caminhos não autorizados, categorias de edição não
suportadas, edições coordenadas incompletas e inconsistências de símbolos entre
arquivos são rejeitados; edições de build, de dependências, de interface
pública e de arquivos novos exigem categorias explícitas. Verificações
semânticas rejeitam formas de falha conhecidas: alterar um valor armazenado,
transmitido ou comparado com uma constante sem manter a forma antiga para os
dados existentes; deixar inalterado um leitor de um valor armazenado alterado;
remover um texto cifrado que o app distribui; substituir um algoritmo exigido
por um protocolo; ou deixar o valor reportado intacto. Cada rejeição retorna ao
gerador com o motivo; quando nenhuma mudança dentro do app consegue satisfazê-la
(texto cifrado sob uma chave distribuída), o achado vai para os mantenedores
pelo encaminhamento descrito abaixo. Verificações sobre um comportamento que o
mantenedor pode ter escolhido (remover ou esvaziar uma configuração ou opção do
usuário, sobrescrever ou recusar um valor fornecido pelo chamador) não recusam o
patch: viram notas mostradas ao autor antes da aprovação, e a revisão de
comportamento ainda conta uma opção esvaziada como regressão. Uma
renomeação deve atualizar todas as referências no mesmo patch, ou é gerada
novamente com a falha explicada; autorizá-la ou não é uma nota no pull request,
não um bloqueio.

Resta uma confiança residual: o índice não é uma análise JVM completa, uma chamada resolvida
pode depender de despacho dinâmico, reflexão ou código gerado, e a documentação
descreve o contrato de uma API, não o provedor implantado. Suposições não
resolvidas permanecem visíveis como \emph{suposições de revisão} para um
mantenedor ou teste, distintas dos \emph{bloqueadores de geração} (o código não
teria sustentação) e dos \emph{bloqueadores de publicação} (construível
localmente, mas não apresentado como pronto).

\subsection{Revisão de Comportamento e Encaminhamento}

Uma regressão exige um cenário concreto alcançável pelo próprio código da
aplicação. Remover algo obsoleto ou inseguro oferecido a um par ou aceito dele
é endurecimento, a menos que as evidências mostrem quem depende disso: o
contrato da função, uma opção no código, outro código que lida com esse par, ou
os mantenedores. A revisão atribui cada regressão ao plano (nenhum rascunho que
o siga consegue evitar o cenário) ou ao rascunho. Um deslize do rascunho gera
um novo rascunho sobre o mesmo plano, com o cenário como feedback, regenerando
os testes cujo arquivo de produção esse feedback altera. Após o último rascunho
do plano, ou após duas regressões consecutivas causadas pelo plano, a unidade
recebe um novo plano com os cenários como restrições---no máximo uma vez, por
qualquer motivo; se o novo plano não for adotado, os rascunhos restantes do
primeiro plano ainda são executados. A revisão e as verificações
determinísticas também rejeitam um patch em que a chamada governada continua
recebendo o mesmo material (uma constante
renomeada, um alias ou uma factory em torno do mesmo objeto, uma chave
reconstruída a partir de seus próprios bytes ou cifrada novamente sob a mesma
chave); se restarem apenas rejeições desse tipo, o achado é encerrado como
violação sem impacto de segurança. Uma mudança que altera o valor, como
\texttt{SSL} para \texttt{TLS}, ainda se torna uma proposta.

Achados sem patch são encaminhados de acordo com o impacto de segurança
estabelecido. Uma violação confirmada interrompida antes de qualquer rascunho
(sem correção segura, nenhuma edição encontrada, um contrato externo) torna-se
uma issue somente quando seu impacto foi confirmado; caso contrário, é
registrada como sem correção segura ou como bloqueada por um contrato externo,
com o respectivo bloqueador. Código que já usa a API conforme documentado é
encerrado como violação sem impacto de segurança somente com um impacto não
relacionado à segurança. Uma geração que falhou mantém sua issue, a menos que
todo bloqueador seja uma verificação dos próprios rascunhos do pipeline (uma
renomeação feita pela metade, um símbolo inventado, o valor reportado deixado
intacto) ou, quando a adjudicação não deixou pergunta ao mantenedor, da forma
deles (um rascunho sem edições). Esse achado termina como correção nunca
produzida, sem issue para ele, pois nenhum mantenedor pode responder por um
rascunho que o pipeline recusou; o mesmo vale, na mesma condição, para uma
correção que a adjudicação apoiou e que parou antes de qualquer rascunho apenas
porque a análise não revisou arquivos que o pipeline calculou como parte do
escopo. Cada achado
termina em exatamente um desfecho: pull request pronto (ou um já aberto); patch
sem pull request (não aplicado, sem build ou retido); não aplicável; correção
recusada antes de qualquer rascunho ou nunca produzida; issue para a aplicação
ou para uma biblioteca embarcada; violação confirmada sem correção segura;
violação sem impacto de segurança; contrato externo; decisão do mantenedor;
falso positivo reconhecido; evidência insuficiente; ou não analisado.

\subsection{Memória entre Execuções}

As respostas dos mantenedores, atualizadas a partir do GitHub, chegam à
adjudicação, ao planejamento e à revisão de comportamento para achados da mesma
especificação, de modo que uma premissa rejeitada não seja proposta novamente e
um achado cujo pull request foi fechado não seja republicado. Tentativas
anteriores são lidas apenas quando chegaram a um patch, para que a amostragem
de uma execução não se torne evidência da execução seguinte; fatos confirmados
por um sinal determinístico (um import aceito por um build bem-sucedido) são
dados e nunca alteram código ou prompts. Para fins de avaliação, uma execução
pode começar sem histórico de tentativas.

\subsection{Aplicação na Ponta do Branch e Build}

A análise lê a revisão do APK, mas um pull request tem como alvo a ponta atual
do branch base. Cada solução aceita é aplicada atomicamente em um worktree Git
dedicado na ponta, nunca no clone base, com um caminho distinto por tentativa e
um registro do commit base, das edições e do motivo de falha. As edições
literais são tentadas primeiro, acompanhando arquivos renomeados, e depois o
merge de três vias do Git a partir da revisão analisada. Somente se ambos
falharem uma etapa estreita do modelo porta o patch, recebendo o patch
revisado, as mudanças upstream e os arquivos da ponta: ela reancora a correção,
declara que a ponta não contém mais a violação, ou recusa o porte. Um porte
deve passar pelas verificações semânticas e de ancoragem; as linhas novas
passam outra vez pela revisão de comportamento, e um porte rejeitado
retorna mais uma vez com o motivo; um porte recusado pelas
verificações enquanto a violação persiste torna-se candidato a issue. Uma
alegação de que a violação já foi corrigida só é aceita quando uma verificação
determinística confirma que a instrução reportada não está em nenhum arquivo da
ponta, ou que todos os blocos corrigidos estão presentes e o código substituído
não; caso contrário, trata-se de um porte recusado. Um achado é
\emph{não aplicável}, sem pull request, quando a ponta já contém a correção, o
arquivo saiu do branch, ou nenhum porte se aplica e o código
substituído existe apenas na revisão analisada---não resta nada a corrigir tal
como foi revisado, o que não equivale a uma violação desaparecida. O resumo é
reescrito após um porte.

O build executa a tarefa Gradle selecionada, normalmente uma montagem de debug
sem minificação, em um JDK escolhido a partir do wrapper e da versão de Java
declarada, registrando comandos, código de retorno, APKs, logs e diagnósticos.
Falhas de SDK, JDK, Gradle, rede ou submódulos são falhas de infraestrutura;
uma falha compartilhada pela ponta sem modificações é do projeto, mas nunca
conta como aprovação. Passagens mais estreitas (apenas o módulo modificado,
apenas o código de produção, ou compilado na revisão analisada porque a ponta
não compila) mantêm seus próprios status e nunca contam como build completo. Um
patch em um source set que o APK não compila não é publicado. Uma etapa
estreita de reparo de compilação pode corrigir novos diagnósticos nos arquivos
alterados; ela não pode reabrir a decisão de segurança, e edições fora do plano
são revisadas novamente. Onde existe uma regra estática para a especificação,
ela é reexecutada antes e depois do patch como medição. Essas verificações
mostram que o patch se aplica e compila, não que ele remove a violação em tempo
de execução ou preserva o comportamento; se a violação ainda aparece é
verificado pela validação dinâmica dos pull requests publicados
(Seção~\ref{subsec:dynval}).

\subsection{Proposta e Publicação}

O texto é escrito para o mantenedor, não como um relatório do pipeline. Um
modelo o redige a partir do diff final e dos fatos da revisão em uma forma
fixa: um título que nomeia a mudança, poucos parágrafos simples sobre o que o
código fazia, o que mudou e se o build passa, e um fecho fixo que nomeia o
RVSec como analisador e apresenta o autor. Quando o repositório tem um template
de pull request, o texto o preenche, mantendo seus títulos e checklists e
marcando apenas caixas verificadas, como um build bem-sucedido, nunca um teste
em dispositivo, já que o RVSec executou o APK original. Verificações
determinísticas então prendem o rascunho a essa forma e ao diff: um título
genérico como ``Fix \texttt{<init>}'', identificadores de evidência, nomes de
bytecode compilado, um resumo com mais de 250 palavras, termos do vocabulário
próprio do pipeline ou uma frase que descreve uma edição
ausente do diff final devolvem o texto para reescrita com o problema indicado,
e a publicação só é retida quando a reescrita ainda falha. O objetivo é um
texto nos termos do projeto, não um que esconda sua origem: o uso de IA é
declarado quando um template o pergunta. Um \emph{pull request pronto} é um rascunho com commit
cujo build passou: um candidato à revisão humana, não uma publicação.

Nada é publicado automaticamente: o operador aprova cada achado por seu
identificador após uma pré-visualização. O \system{} então verifica novamente o
repositório ao vivo e, por padrão, recusa um achado cujo pull request anterior
possa estar aberto ou tenha sido fechado sem merge, uma mudança que outro de
seus pull requests faz ou com a qual se sobrepõe, um diff que o branch
padrão já contém ou com o qual não corresponde mais, um diff apenas de testes,
uma revisão feita sob um contrato mais antigo, ou um repositório que não aceita
contribuições; cada recusa indica a opção que a anula. Uma política sobre
contribuições de IA é informada ao operador, que então escreve o texto; ela não
bloqueia. Uma issue é recusada quando seu código saiu do branch padrão. O
\system{} nunca assina o commit de outra pessoa; comentários de revisão são
respondidos revisando o pull request, novamente após aprovação. A criação tem
como alvo o GitHub; em outras plataformas, a análise e o patch validado ainda
são produzidos. Os registros de prompts guardam requisições, respostas,
versões, tokens e latência, nunca credenciais.

\subsection{Entradas e Execução}

O \system{} não depende de um único detector. Um achado precisa de quatro campos
(o projeto, a classe, o método e a especificação violada) e pode trazer também a
mensagem do detector, o repositório e o commit; os achados são lidos de saídas
CSV, JSON, SARIF ou Semgrep, ou produzidos por \texttt{cognalys scan}. Para o
RVSec, \texttt{cognalys fdroid-build} compila APKs sem minificação a partir das
receitas do F-Droid, \texttt{cognalys github-build} faz o mesmo a partir do
branch padrão de um aplicativo Android no GitHub sem receita, e
\texttt{cognalys rvsec-run} executa o RVSec com os exploradores escolhidos e
importa seus relatos; a instrumentação e as
especificações monitoradas ficam no RVSec.

Uma execução é iniciada com \texttt{cognalys run} e um arquivo de configuração
curto. Outros comandos, também disponíveis como API Python
(\texttt{cognalys.api}), listam as etapas de uma execução, resumem desfechos,
informam o custo, mostram ou aplicam patches, publicam propostas aprovadas e
atualizam as respostas dos mantenedores.

As chamadas de modelo vão para a OpenAI ou a Anthropic, ou para um modelo local
via Ollama, que roda sem pesquisa na web. As chamadas à OpenAI usam o nível Flex,
cobrado pelas tarifas da API Batch~\cite{openai2026flex}; uma chamada que o Flex não consegue atender passa
para o nível padrão, e o rastro registra qual nível respondeu. Duas
salvaguardas limitam uma execução: ela para de chamar modelos ao atingir um
teto de custo configurado, e uma etapa de modelo sem chamada nem linha de log
por 30 minutos é interrompida e retomada uma vez a partir do seu cache de
respostas. As respostas bem-sucedidas do modelo ficam em cache entre execuções,
sob um hash da requisição inteira, separadamente do corpus de pesquisa.

\subsection{Um Exemplo Completo}

O RVSec relatou que \texttt{CryptoManager.getSecretKey}, no Traffic Light
(version code 38), executava uma sequência de chamadas que a
\texttt{KeyStoreSpec} não aceita (a regra rejeita um \texttt{getEntry} não seguido de
um \texttt{getKey}~\cite{costa2026effectiveness}): lia a chave com \texttt{KeyStore.getEntry} e
retirava a chave secreta da entrada retornada. O \system{} mapeou o método para
o código-fonte Kotlin, confirmou a violação e sua causa, e julgou seu impacto
como não relacionado à segurança: a chave é apenas lida, e o app não a cria nem
a protege de outro modo. Como uma mudança de uma linha levava o código a uma
sequência aceita, a unidade seguiu para o reparo. O rascunho trocou a leitura e
manteve o cache e o fallback que cria uma chave ausente; passou pela revisão de
comportamento, foi aplicado na ponta do branch e compilou.

{\scriptsize
\begin{verbatim}
- (keyStore.getEntry(key, null)
-     as? KeyStore.SecretKeyEntry)?.secretKey
+ (keyStore.getKey(key, null) as? SecretKey)
      ?: createKey(key)
\end{verbatim}}

O pull request (\#238), com cerca de 60 palavras em parágrafos simples e
nomeando o RVSec como analisador, foi aberto em 20 de outubro de 2026 e
mergeado no dia seguinte, com uma resposta de uma palavra. Reexecutado sobre o
APK com o patch, o RVSec não relatou mais a violação no método em nenhuma de
três execuções (Seção~\ref{subsec:rq7}). O relato desaparece porque o monitor
passa a ver de onde veio a chave, não porque o app ficou mais seguro: trata-se
de uma violação confirmada sem impacto de segurança, que ainda assim vira pull
request quando sua correção é segura.

\section{Configuração do Estudo}
\label{sec:settings}

O estudo tem quatro partes, com dados e datas distintos. (i)~Um estudo de caso
retrospectivo dos 34 achados distintos do retrato consolidado do
experimento de 19 de outubro de 2026, montado em 22 de outubro, quando o pipeline
ainda estava mudando (\researchq{1}, \researchq{2}, \researchq{4}, \researchq{5}).
(ii)~Um levantamento dos 49 pull requests e 10 issues submetidos para achados do
RVSec, coletados do GitHub em 18 de novembro de 2026 (\researchq{3};
Seção~\ref{subsec:selection}). (iii)~Uma comparação pré-registrada, com juízes
cegos, entre uma configuração registrada do {\system} e agentes de programação em dois
blocos, três braços que recebem o achado e o repositório e dois instruídos
a seguir o processo do {\system}, sobre 20 unidades de análise
(\researchq{6}; Seções~\ref{subsec:frozen} e~\ref{subsec:protocol}).
(iv)~Uma validação dinâmica pré-registrada que reexecuta o RVSec sobre o APK original e
o APK com patch dos pull requests publicados (\researchq{7};
Seção~\ref{subsec:dynval}). As execuções anteriores a um registro, e todo
desvio, são declarados; eventos posteriores não reescrevem um resultado anterior de build ou de
adjudicação. As quatro partes não medem um único sistema em um único momento.
O estudo de caso descreve o pipeline como estava em outubro de 2026, com um
modelo anterior e várias revisões da biblioteca; os pull requests foram
produzidos pelas versões em uso quando cada um foi aberto, de agosto a novembro
de 2026; a comparação executa uma configuração fixa com o modelo atual; e a
validação dinâmica reexecuta o RVSec sobre os pull requests publicados,
qualquer que seja a versão que os produziu. Os números são comparados dentro de
uma parte, não entre partes.

\subsection{Objetivos e Questões de Pesquisa}
\label{subsec:rqs}

O objetivo de implementação é complementar a detecção em tempo de execução com
classificação no nível do código-fonte, abstenção e reparo revisável por meio de uma
biblioteca reutilizável. O objetivo empírico é descrever o que ela decidiu,
construiu e gastou e como os mantenedores responderam; verificar com o RVSec se os
patches publicados removem a violação; e compará-la com
agentes de programação de propósito geral, com e sem o seu processo como instruções.

\begin{description}
  \item[\researchq{1}:] \emph{Que resultados de análise e de reparabilidade
  o {\system} produz para os achados de uso indevido de API selecionados?}
  Uma decisão de proposta não implica um patch gerado ou validado.

\item[\researchq{2}:] \emph{Em que medida os patches gerados podem ser
  aplicados e compilados nos projetos-alvo?}
  Relatamos a aplicação, o escopo do build (APK ou módulo alterado) e a atribuição
  das falhas; trata-se de compilabilidade, não de correção comportamental ou de segurança.

  \item[\researchq{3}:] \emph{Como os mantenedores respondem às propostas de reparo
  submetidas, e que razões apresentam para suas decisões?}

  \item[\researchq{4}:] \emph{Quais relatos do detector a análise
  reconhece como falsos positivos ou como violações sem impacto de segurança,
  e eles se concentram em especificações particulares?}

  \item[\researchq{5}:] \emph{Quanto uso de modelo a análise
  requer, e como ele se distribui entre os resultados?}

\item[\researchq{6}:] \emph{(a) Um agente de programação de propósito geral basta quando
  recebe o achado e o repositório? (b) Executar o processo como um
  pipeline acrescenta algo em relação a um agente instruído a seguir o mesmo processo?
  (c) Quanto variam as entregas do {\system} entre execuções?}
  Uma configuração de agente usa o modelo do {\system}, de modo que uma diferença entre os
  dois não pode ser atribuída ao modelo.

  \item[\researchq{7}:] \emph{Quando o RVSec é executado novamente sobre o APK com patch de um
  pull request publicado, a violação relatada ainda aparece no
  método alterado pelo patch?}
  Um build ou um merge não mostra que o uso indevido observado desapareceu.
\end{description}

A identidade de um achado é o identificador do APK, a classe, o método e a especificação; a
identidade de um pull request é o repositório e o número do pull request. Tentativas repetidas não são casos
independentes; especificações diferentes em um mesmo método permanecem achados distintos, mesmo
quando um único patch trata de ambas. O estado do pull request e o estado do build são dimensões
separadas. O estudo de caso e o levantamento são exploratórios: nenhuma hipótese foi
fixada de antemão, e o pipeline mudou nesse meio-tempo (cada mudança está registrada
ao lado da medição que a motivou). A comparação e a validação
dinâmica seguem procedimentos e regras de análise escritos e datados
antes de suas execuções; os registros da comparação também guardam os hashes de
seus prompts, scripts e gabaritos. Cada registro informa quais execuções
o precederam, e toda regra alterada depois de os resultados serem vistos é declarada como
desvio. Nenhum teste estatístico é aplicado: os resultados são contagens descritivas e
explicações de casos (Tabela~\ref{tab:rq-data}).

\begin{table}[H]
\centering
\caption{Dados por trás de cada questão de pesquisa.}
\label{tab:rq-data}
\footnotesize
\begin{tabular}{p{0.9cm}p{3.7cm}p{2.6cm}}
\toprule
QP & Dados & Unidades e data \\
\midrule
QP1, QP4 & Estudo de caso: retrato consolidado do experimento de 19 de outubro & 34 achados, 22/10 \\
QP2 & Estudo de caso: correções geradas, aplicação na ponta e builds & 11 correções, 7 com build \\
QP3 & Levantamento no GitHub de pull requests e issues para achados do RVSec, três origens (Tabela~\ref{tab:origin}) & 49 PRs, 10 issues, 18/11 \\
QP5 & Rastros de prompts das execuções do estudo de caso & 34 achados, 11 execuções \\
QP6 & Comparação com agentes de programação (registrada de antemão) & 20 unidades (24 achados), 4 execuções por unidade \\
QP7 & RVSec sobre APKs originais e com patch (registrado de antemão) & 45 de 49 PRs, 14--19/11 \\
\bottomrule
\end{tabular}
\end{table}

\subsection{Seleção de Casos}
\label{subsec:selection}

\subsubsection{Estudo de caso datado e levantamento de pull requests}

O estudo de caso reúne 34 achados distintos de 23 APKs, cada um com o
repositório e o commit de origem do seu APK. Os achados vêm da campanha do grupo do
RVSec sobre 163 aplicações do F-Droid~\cite{costa2026effectiveness}. Nos dados da campanha, dos 109 usos
selecionados para acompanhamento no código-fonte, 72 estavam em
código de biblioteca ausente do repositório no commit analisado, três em uma cópia
vendorizada de uma biblioteca e 34 no código da própria aplicação (31 usos
indevidos que a campanha estabeleceu e três vistos em uma única execução, que
ela deixou indecisos). O estudo de caso combina uma execução base e reexecuções
direcionadas posteriores, mantendo o resultado mais recente de cada achado; assim, contém 34
achados, e não 34 execuções independentes, e descreve o que a
versão de outubro de 2026 decidiu, construiu e gastou, não o que a versão atual
decidiria; a \researchq{6} mede a configuração fixa da
Seção~\ref{subsec:frozen}. Essas execuções mantiveram o cache de respostas, de
modo que uma requisição ao modelo idêntica a uma já respondida foi atendida por
ele em vez de ser amostrada de novo (no caso da pesquisa na web, apenas dentro
de uma semana).

Os estados e as respostas dos pull requests foram coletados pela API do GitHub em
18 de novembro de 2026. A inclusão exige um pull request aberto, uma mudança incorporada ou fechamento
por um mantenedor; tentativas substituídas que foram fechadas pelo autor ou por um bot são
excluídas. O BombusMod \#316 é uma incorporação manual, não um merge no GitHub.

Um pull request conta como RVSec quando seu achado vem de um relatório do RVSec,
rastreado até a execução que produziu o pull request. Os pull requests e as
issues têm três origens (Tabela~\ref{tab:origin}). Os dados da campanha do grupo do
RVSec~\cite{costa2026effectiveness} fornecem os 34 achados do estudo de caso;
oito pull requests e três issues vêm de execuções sobre os seus dados depois de
16 de setembro de 2026. Nossas próprias execuções de 8 de agosto a 16 de
setembro de 2026, sobre relatórios do RVSec anteriores à campanha, produziram
15 pull requests e cinco issues; oito desses pull requests tratam de achados do
estudo de caso, que a campanha relatou de novo, e se ligam a eles pela
identidade do achado. Nossas filas do RVSec produziram os outros 26 pull
requests (abertos de 16 de outubro a 7 de novembro) e duas issues: cada
aplicação foi executada uma vez sob o RVSec, com Monkey e APE, por 600
segundos, com as especificações criptográficas e de TLS, e apenas achados no
código da própria aplicação foram analisados. Os candidatos das filas foram
aplicações do F-Droid e, depois, aplicações Android do GitHub sem receita no
F-Droid, preferindo aquelas cujo código criptográfico executa sem conta; mais
tarde, exigiu-se também que seus repositórios estivessem ativos e tivessem
incorporado, nos doze meses anteriores, um pull request de um colaborador
externo.

\begin{table}[H]
\centering
\caption{Origem dos pull requests, das issues, das unidades da comparação e
dos casos da validação dinâmica. Pull requests e issues são atribuídos pela
execução que os produziu; as unidades da comparação, pela fonte do seu achado.
Estados em 18 de novembro de 2026: incorporados/fechados/abertos. QP7: casos
executados/rotulados como corrigidos.}
\label{tab:origin}
\footnotesize
\setlength{\tabcolsep}{3pt}
\begin{tabular}{p{3.0cm}rcrrc}
\toprule
Origem & PRs & Inc./Fech./Abertos & Issues & QP6 & QP7 \\
\midrule
Nossas execuções sobre relatórios anteriores do RVSec (8/8--16/9) & 15 & 6/4/5 & 5 & 0 & 12/6 \\
Campanha do grupo do RVSec (após 16/9) & 8 & 4/1/3 & 3 & 10 & 7/5 \\
Nossas filas do RVSec (16/10--7/11) & 26 & 11/5/10 & 2 & 10 & 26/22 \\
\midrule
Total & 49 & 21/10/18 & 10 & 20 & 45/33 \\
\bottomrule
\end{tabular}
\end{table}

A Tabela~\ref{tab:sample} apresenta o inventário.

\begin{table}[H]
\centering
\caption{Inventário de casos e sobreposição (as contagens de PRs e de achados usam unidades diferentes).}
\label{tab:sample}
\footnotesize
\begin{tabular}{p{6.1cm}r}
\toprule
Componente do inventário & $n$ \\
\midrule
Achados consolidados & 34 \\
Registros de PRs publicados elegíveis & 49 \\
Registros de PRs com identidade de achado completa & 48 \\
Achados sobrepostos, mantidos uma única vez & 15 \\
Achados adicionais de PRs totalmente vinculados & 32 \\
\quad das nossas filas do RVSec & 26 \\
\quad das nossas execuções anteriores & 6 \\
União de achados totalmente vinculados & 66 \\
Identificadores de APK distintos nessa união & 53 \\
Registros de PRs com vínculos de origem incompletos & 1 \\
\bottomrule
\end{tabular}
\end{table}

A correspondência usa a tupla completa do achado e a URL publicada, a identidade do retorno
ou os metadados de origem, nunca apenas o pertencimento ao repositório; vínculos não resolvidos permanecem
visíveis em vez de receberem identidades inventadas. Os 48 registros com
identidade completa correspondem a 47 achados, 15 deles sobrepostos ao estudo de caso,
de modo que a união contém 66 achados de 53 rótulos de APK. O HTTP Shortcuts \#576,
sem um vínculo completo de APK/classe/método/especificação, permanece no inventário
de PRs, mas fora desses denominadores. A seleção é intencional (casos
efetivamente analisados ou submetidos, incluindo retorno negativo), não uma
amostra representativa. Quarenta e seis dos 49 pull requests tratam de
especificações criptográficas (a do HTTP Shortcuts \#576, \texttt{SSLContextSpec},
foi recuperada de uma fonte registrada para a \researchq{7}); três (BombusMod,
wwwjdic e ntfy), das nossas execuções anteriores, tratam de especificações de
manipulação de recursos e de sincronização que nossas filas não habilitam. Quatorze pull requests para candidatos identificados
manualmente, por pesquisa na web e por regras do Semgrep, submetidos durante o
desenvolvimento, são relatados à parte, como extra, e ficam fora de todas as contagens: em
18 de novembro de 2026, cinco haviam sido incorporados (aiHub \#279, HomeApp \#114,
Auto-Shine \#12, SlimBook \#3 e aaf-easydiary \#253), dois estavam abertos e
os mantenedores fecharam sete.

\subsubsection{Procedência de APKs e revisões}

O RVSec observa um APK, mas o \system{} corrige um commit, por isso cada achado
registra o repositório e o commit do seu APK. Dos 26 dos 53 APKs cujo
binário ainda está disponível, o manifesto coincidiu em 22; os outros quatro são
builds de depuração cujo pacote traz um sufixo de variante (\texttt{.dev},
\texttt{.staging}, \texttt{.debug}), dois deles também com um código de versão diferente.
Para os outros 27, a correspondência é presumida
(Seção~\ref{subsec:threats}). As aplicações das filas vindas do GitHub sem
receita no F-Droid foram construídas a partir da ponta do seu branch padrão;
três dos 49 pull requests vêm delas (TriliumDroid \#177, meshcore-mobile-app
\#10 e CometOTP \#114). Todos os repositórios estão no GitHub, cuja API é usada
tanto para publicar quanto para coletar as respostas. Quando o branch base avançou, o
patch é levado adiante antes de se abrir um pull request
(Seção~\ref{sec:pipeline}).

\subsection{Procedimentos de Coleta de Dados}
\label{subsec:collection}

\begin{table}[H]
\centering
\caption{Principais artefatos coletados.}
\label{tab:artifacts}
\footnotesize
\begin{tabular}{p{2.3cm}p{5.0cm}}
\toprule
Artefato & Informação usada no estudo \\
\midrule
Método extraído & Metadados do achado, código-fonte exato, imports, pontos de chamada,
contexto da especificação, orientações do repositório \\
Pacote de evidências & IDs estáveis de evidência, procedência, confiança, relações,
ambiguidades, exclusões, completude, hashes \\
Adjudicação & Estados de causa e de impacto, decisão, bloqueios, alegações de suporte,
testes exigidos \\
Plano de patch/solução & Arquivos e edições autorizados, invariantes, estado da revisão,
substituições exatas \\
Relatório de execução & Worktree/commit base, estado da aplicação, arquivos alterados,
revisão e erros \\
Relatório de build & Tarefa, estado, código de retorno, duração, diagnósticos, logs e
atribuição de falhas \\
Rastro de prompt & Etapa, requisição/resposta, versões, tokens, latência, cache e
identificadores de requisição \\
Cartão julgado & Tipo de entrega, título, texto, diff, indicador de build; notas e observações dos juízes \\
Reexecução do RVSec & APK original e com patch, violações e cobertura no método alterado pelo patch \\
\bottomrule
\end{tabular}
\end{table}

Cada achado guarda os artefatos da Tabela~\ref{tab:artifacts}. Para a
\researchq{3}, registramos, para os 49 pull requests elegíveis e as 10 issues,
abertos entre 8 de agosto e 7 de novembro de 2026 (as issues até 5 de novembro), seu
estado, o merge, quem os fechou
e se alguém além do autor ou de um bot comentou, revisou ou os
fechou (uma resposta de mantenedor); as issues são consultas sobre se um pull
request seria bem-vindo ou relatos de achados sem correção segura. As 23 issues
abertas em 4 de agosto, um teste com um modelo local, ficam de fora
(Seção~\ref{subsec:ethics}). Para a
\researchq{5}, cada chamada rastreada registra seu modelo e os tokens de entrada, de entrada em cache e
de saída, e é atribuída ao seu achado; as chamadas enviadas a um provedor são precificadas pelos preços
de tabela e somadas por achado ao longo das execuções do retrato, de 19 e 20 de outubro de 2026, excluindo chamadas feitas para nenhum
achado e requisições respondidas pelo cache de respostas, que não têm custo. Para a \researchq{6}, cada execução de agente guarda sua entrega, diff, log de
build, log de eventos e tokens, além da fase de porte no caso de um pull request; a
chave cartão-braço é mantida separada dos juízes. Para a \researchq{7}, cada caso
guarda os dois APKs, o registro do build com patch e todos os relatórios do RVSec.

\subsection{Procedimentos de Análise de Dados}
\label{subsec:analysis}

Mantemos o resultado mais recente de cada achado sem agregar tentativas, e o estado do pull request
e o retorno separados desse retrato; o inventário legível por máquina guarda
caminhos de origem, base da correspondência, URLs dos PRs, disponibilidade do binário, pacote do manifesto
e versão. Uma proposta com build válido não é necessariamente correta, segura,
aceita ou livre do relato do detector (verificado para os pull
requests publicados na \researchq{7}, apenas nos caminhos que os exploradores alcançam), e um pull request
fechado não é necessariamente um patch errado (comportamento intencional, política, uma
mudança já aplicada ou incorporação manual). Por isso evitamos uma taxa de sucesso
agregada. Os desfechos de \researchq{1} e \researchq{4} são as decisões
registradas pelo pipeline; quando o texto chama uma decisão de correta ou não,
essa é a leitura do autor sobre as evidências da execução. As razões de
fechamento da Tabela~\ref{tab:closures} foram codificadas pelo autor a partir das
palavras de cada mantenedor, uma razão por pull request, sem um segundo
codificador.

\subsection{Configuração Fixa e Unidades da Comparação}
\label{subsec:frozen}

A comparação executa o {\system} em uma configuração fixa: uma única versão da biblioteca
(quatro de suas 80 execuções, que haviam falhado por defeitos de
implementação, foram reexecutadas depois da correção desses defeitos), \texttt{gpt-6.1-terra} em todas as etapas, com esforço de raciocínio
alto e limite de saída de 32.768 tokens, no nível Flex, com todo patch
passando por build e um teto de US\$~2 por execução. Uma execução só conta quando o manifesto do
experimento (modelo, prompts, schemas, políticas, especificações, repositório de aprendizado,
corpus de pesquisa, commit) coincide. Aplicam-se três controles de vazamento: o próprio
repositório da unidade é excluído (nem as respostas dos seus mantenedores nem o que execuções
anteriores aprenderam ali chegam a qualquer prompt, enquanto as respostas de outros repositórios à
mesma regra permanecem, e os agentes do segundo bloco recebem as mesmas); o
repositório de aprendizado guarda imports verificados, mas nenhuma tentativa anterior; e nenhuma resposta
de uma execução anterior é reutilizada, exceto a pesquisa na web compartilhada, no nível da regra, do
corpus registrado, sendo o caso da própria unidade sempre pesquisado de novo; dentro de
uma execução, uma requisição que a análise de uma unidade repete sem alteração, como
em um replanejamento, reutiliza a primeira resposta.

O modelo foi escolhido por uma verificação descritiva: oito unidades (dez achados de
seis aplicações) foram reexecutadas com \texttt{gpt-5.6-terra}, usado até
2 de novembro de 2026, e com \texttt{gpt-6.1-terra}
sob as mesmas configurações. Ambos fizeram 83 chamadas, com custo de US\$~6,31 e US\$~2,43
aos preços de tabela (US\$~2,71 para o \texttt{gpt-6.1-terra} com seu preço de escrita em cache
e a taxa de busca); concordaram em quatro das oito unidades. A amostra é pequena demais para comparar a qualidade das decisões;
o custo menor sustentou a troca.

As 20 unidades de análise foram sorteadas em 13 de novembro de 2026 com uma semente registrada.
Dez vêm de pull requests publicados cuja execução original terminou com um
rascunho com commit e build bem-sucedido, estratificados pelo estado nessa data: três incorporados, cinco
abertos, dois fechados pelo mantenedor. Dez vêm dos 45 achados distintos do estudo
de caso (34) e de uma comparação exploratória anterior com agentes (12, um deles
em comum), que o pipeline analisa novamente, excluindo métodos que também
aparecem na população de pull requests dos dez primeiros: nove sem
pull request e um cujo pull request o mantenedor fechou. Por origem
(Tabela~\ref{tab:origin}), dez unidades vêm da campanha do grupo do RVSec (duas
delas entre as unidades de pull request) e dez das nossas filas (oito). Uma unidade é um achado, ou
vários achados que o pipeline reúne em um único patch
(Seção~\ref{sec:pipeline}): quatro unidades reúnem dois achados cada, duas dentro
de um mesmo método e duas em dois métodos de uma mesma classe.

\subsection{Comparação com Agentes de Programação}
\label{subsec:protocol}

\emph{Braços.} Cinco braços de agentes, em dois blocos, são comparados com o {\system} (braço
C), em três configurações: o Codex CLI com \texttt{gpt-6.1-terra}
(o modelo do {\system}), o Claude Code com Claude Sonnet 5.5 e o Codex CLI
com \texttt{gpt-6.1-sol}. No primeiro bloco (T, S, O), cada agente recebe
o achado e o repositório e faz uma única entrega, julgada no commit analisado,
sem porte; no segundo (TP, SP2), as duas primeiras configurações
recebem também o processo do {\system} como instruções e portam seus pull requests
para a ponta do branch. Todos executam com esforço alto, quatro vezes por unidade. T, TP e C
compartilham o modelo, de modo que suas diferenças refletem as instruções e o local onde o
processo é executado (uma única sessão, ou etapas separadas com verificações determinísticas
entre elas). Uma execução parcial do {\system} com Claude Sonnet 5.5 (17 das 40
execuções planejadas) não é comparada; seus cartões entram apenas na concordância
entre os juízes e na sondagem de cegamento do primeiro bloco.
Um terceiro braço do segundo bloco, com \texttt{gpt-6.1-sol}, foi interrompido após 25 de
suas 80 execuções, antes que qualquer uma fosse julgada, quando a cota da assinatura se esgotou,
e não é comparado (um desvio declarado).

\emph{Fase 1.} Cada execução parte de um checkout limpo do commit analisado,
sem remoto, credenciais ou memória, com limite de 60 minutos, busca na web
e o comando de build do {\system}. O prompt fornece as linhas do relatório do RVSec, pede
exatamente uma entrega (pull request, issue ou nada a enviar) e enuncia
as regras de texto do {\system}. O prompt do segundo bloco acrescenta a regra publicada
de cada relatório, as respostas de mantenedores de outros repositórios a ela e as entradas que o {\system}
lê, e descreve cada etapa em ordem (evidências, pesquisa, decisão,
planejamento, rascunho, verificações, revisão de comportamento, novo rascunho e replanejamento, build,
texto). A primeira versão do braço com Claude Sonnet, o SP, pedia, como o TP,
anotações etapa a etapa; as salvaguardas do modelo no Claude Code interromperam a
maioria de suas execuções, por isso o SP2 usa o prompt sem esse pedido. Um prompt
não consegue reproduzir contextos separados por etapa, validadores e
novas tentativas impostos por código, nem o corpus de pesquisa.

\emph{Fase 2.} Assim como faz o {\system} (Seção~\ref{sec:pipeline}), os pull requests
do segundo bloco vão para a ponta única da unidade, obtida das execuções do {\system}:
primeiro um merge de três vias, com build sem o agente quando não há conflito; havendo conflito, o
mesmo agente, com o mesmo modelo e esforço, informa o patch como portado, já corrigido ou
abandonado. ``Já corrigido'' só se mantém quando uma verificação determinística encontra as
linhas adicionadas presentes e as linhas removidas ausentes na ponta; caso contrário, ele é
abandonado. Os tokens de ambas as fases entram no custo.

\emph{Cartões e juízes.} Cada entrega se torna um cartão (tipo, título, texto,
diff em relação ao commit analisado, indicador de build), despojado de assinaturas,
identificadores, nomes de ferramentas e vocabulário do pipeline, o que é garantido por um lint. Dois
juízes, Claude Opus 5.5 (principal) e \texttt{gpt-6.1-sol}, com esforço alto,
sem acesso à web e com um checkout somente leitura, pontuam todos os cartões de uma unidade em
ordem aleatória, pois a ordem em que os candidatos são apresentados pode mudar
um veredito~\cite{wang2024fair}. As notas são: patch de 0 (rejeitar) a 3 (mergeável como
está), precisão do texto de 0 a 2 e a justificativa de uma issue ou abstenção
de 0 a 2. A pontuação final é a acordada, ou a do Opus quando divergem. Um
modelo tende a favorecer suas próprias gerações~\cite{panickssery2024selfpref},
e cada juiz é próximo de alguns braços: o GPT-6.1 Sol é o modelo de O e vem do
provedor dos modelos de T, TP e {\system}, enquanto o Claude Opus 5.5 vem do
provedor dos de S e SP2; por isso, a concordância entre os juízes é relatada
ao lado da pontuação final. A pontuação segue uma
referência escrita por unidade (o que um patch aceitável deve mudar e não deve
quebrar), redigida pelos juízes e escrita sem nenhuma resposta de mantenedor.

\emph{Aceitabilidade.} Uma execução é aceitável quando entrega uma issue ou uma
abstenção, ou um pull request com patch de pelo menos 2 e texto de pelo menos 1,
com build na ponta após o porte (primeiro bloco: no commit analisado); uma
falha de build que o código não modificado também apresenta não é atribuída ao patch. Um pull
request cujo porte encontrou a correção já presente na ponta, ou que foi abandonado, nunca
é enviado e é julgado apenas pelo patch (pelo menos 2); essa regra foi definida
depois que os resultados foram vistos e é declarada como desvio.
Uma execução que falhou conta 0. A medida de um braço é a fração de execuções aceitáveis por
unidade, com média sobre as unidades, sob duas leituras: a leitura completa julga o
cartão inteiro, como acima, e a leitura só do patch oculta título e descrição
(aceitável: patch de pelo menos 2 e com build; issues e abstenções não julgadas e
aceitáveis). Para o {\system}, um pull request é qualquer correção gerada, a unidade contada
na \researchq{1}: um rascunho com commit, mesmo um que o pipeline reteria, que
precisa passar no build na ponta, salvo quando a ponta já tinha a correção ou
não tinha mais o código (então é julgado apenas pelo patch). Como issues e abstenções sempre
contam, a medida recompensa a abstenção (um braço que nunca propõe obtém 1),
por isso relatamos também, por braço, os pull requests propostos, julgados aceitáveis e
julgados inaceitáveis (propostos menos aceitáveis), e as unidades com pelo menos um
pull request aceitável. ``Aceitável'' é o veredito dos juízes, não o de um
mantenedor.

\emph{Cegamento.} Uma sessão separada do juiz principal, recebendo um cartão por
vez, tenta identificar seu braço em cartões completos e só de patch, relatando sua
acurácia geral e por palpite. Ela foi executada sobre os cartões do primeiro bloco e do {\system}
e, antes que qualquer pontuação do segundo bloco fosse lida, sobre 10
cartões por braço de TP, SP2 e {\system} em cada visão (um acréscimo declarado).

\emph{Custo e análise.} O custo é uma referência, não um critério: os agentes são
precificados pelos preços de tabela com a taxa de busca na web (os agentes rodaram em assinaturas,
então esse é o custo dos mesmos tokens pela API), e o {\system} pelo que o
nível Flex cobrou e pelo preço de tabela. Com 20 unidades e quatro réplicas, nenhum
teste estatístico é aplicado; relatamos as pontuações finais, a concordância entre os juízes
(kappa de Cohen) e as diferenças observadas. O RVSec não entra nas pontuações, pois apenas
os patches publicados do {\system} foram reexecutados (\researchq{7}).

\subsection{Validação Dinâmica dos Pull Requests Publicados}
\label{subsec:dynval}

O procedimento da \researchq{7} foi registrado em 15 de novembro de 2026, antes que qualquer
caso novo fosse construído ou executado; ele estende um piloto com as seis prévias do retrato que geraram um APK, extensão
decidida quando os resultados do piloto já eram conhecidos. Dos 49 pull requests da
Tabela~\ref{tab:pr-state}, dezesseis são excluídos com motivo registrado:
especificação fora do conjunto criptográfico e TLS que o RVSec verifica (três), ausência de
identidade ou metadados do pipeline registrados, nenhuma linha correspondente do detector no
método relatado, nenhum build que produza o APK analisado, ou um pull request
que trata de outro achado. Os 33 restantes são 28 casos novos e os cinco
casos do piloto cujo patch validado é o do pull request publicado (quatro
corrigidos e o BitBanana, inconclusivo), que o piloto executou três vezes com
as mesmas configurações do RVSec (o piloto usou um JDK fixo; os casos novos o
escolhem como o pipeline faz). O caso PassNotes do piloto verificou outro achado
e outro patch que não os do keepassvault \#347, acrescentado depois como caso
próprio.

O APK original e o APK com patch são construídos a partir do mesmo commit e da mesma tarefa Gradle,
sem minificação; o patch é o diff público do pull request em sua
base de merge, ou o diff do pipeline quando idêntico. O RVSec é executado em cada APK em um
emulador limpo por 600 segundos, com Monkey e APE e suas especificações criptográficas e
de TLS. Cada caso recebe uma repetição, e um caso inconclusivo ou
que não se reproduz recebe uma segunda, se houver tempo. Um build com falha é tentado novamente
uma vez e, depois, relatado como não construído, assim como um patch que não se aplica; o
código da aplicação nunca é editado. Os rótulos seguem uma regra fixada antes das
execuções: \emph{base não reproduz} (fora da contagem) quando o APK original
não mostra violação da especificação no método em nenhuma repetição;
caso contrário, \emph{não corrigido} quando o método com patch relata uma em alguma
repetição, \emph{corrigido} quando ele executa sem nenhuma e
\emph{inconclusivo} quando nunca executa. Os resultados significam ``a violação
relatada não aparece mais no método com patch'', e não ``o RVSec não relata mais
a aplicação''. Casos não concluídos até 18 de novembro de 2026
(19 de novembro para os seis pull requests recuperados e as repetições extras
descritas abaixo) são relatados como não
executados.

\emph{Revisão dos rótulos.} Após as execuções, todo rótulo diferente de
\emph{corrigido} foi revisado com base nas execuções, no diff do pull request e no
código-fonte original, e corrigido quando o rótulo da regra descrevia mal o caso;
o rótulo da regra é mantido ao lado. Quando o método com patch nunca executa, o
caso é \emph{corrigido} se o patch removeu a chamada governada de seu arquivo ou
desativou, por padrão, a condição que protege todas as chamadas do método, e
\emph{não corrigido} se a violação relatada aparece em outro método do mesmo
arquivo. Uma mensagem relatada que desapareceu, deixando no método apenas
mensagens da especificação que o APK original já mostrava, é \emph{corrigido}; um
pull request que altera apenas testes é rotulado como tal. Um caso é \emph{não
verificável} quando a violação relatada não aparece em nenhum dos APKs e sua
mensagem não é emitida por nenhuma execução da validação: o relatório veio de um
conjunto de regras do RVSec anterior ao da campanha.

\emph{Acréscimos após o registro.} Os passos a seguir não faziam parte do
procedimento registrado. O achado do FreeOTP, um dos do estudo de caso, veio da
campanha do grupo do RVSec~\cite{costa2026effectiveness}; no relatório original
da campanha, apenas o DroidBot, o Humanoid~\cite{li2019humanoid} e o
ARES~\cite{romdhana2022ares} alcançaram o método, de modo que seus dois pull
requests são também executados com o DroidBot, um dos exploradores do relatório
original, por 300 segundos e três repetições, ao lado dos resultados de Monkey e
APE (o Humanoid e o ARES não puderam ser executados aqui). Casos inconclusivos e
casos cujo APK original não reproduziu a violação receberam mais uma repetição,
rotulados com todas as repetições juntas: uma terceira para o Read You \#1334,
o Your Local Weather \#209, o InterTune \#52 e o News \#1716, e uma quarta para
o BitBanana, inconclusivo no piloto; no Your Local Weather, o APK original
mostrou a violação apenas na terceira. Depois que os resultados dos
33 já eram conhecidos, 13 pull requests excluídos foram acrescentados: sete cujo
relatório registrado usava um formato de mensagem anterior do RVSec, ao qual
nenhuma linha atual do detector correspondia, rotulados apenas pela especificação
(contando qualquer violação dela no método), e seis cujo APK analisado ou
identidade do achado foi recuperado dos registros do pipeline; doze deles foram
construídos e executados com o RVSec atual nas mesmas configurações. Dos 33 casos
rotulados ao final como corrigidos (Tabela~\ref{tab:dynval}), 27 foram
classificados, 21 deles com uma segunda leitura adversarial, como
alterando o valor ou material governado, alterando a sequência de chamadas ou a
API para a forma exigida, removendo a chamada governada, ou apenas realocando-a
para fora da visão do monitor; essa classificação é descritiva, e os seis
rotulados como corrigidos depois dela não foram classificados
(Seção~\ref{subsec:rq7}). Ao todo, 46 dos 49
pull requests entraram na validação: 45 foram executados, o MetadataRemover \#264
não compila no seu commit base e os três fora do conjunto do RVSec continuam
excluídos.

\section{Resultados}
\label{sec:results}

\researchq{1}, \researchq{2} e \researchq{4} usam o retrato de outubro de
2026, e \researchq{5} o uso de modelo de seus achados: os 34 achados
do experimento de 19 de outubro de 2026, com o resultado da execução mais
recente de cada um no experimento consolidado (19 e 20 de outubro); o
detalhamento dos pull requests por especificação na \researchq{4} usa o
levantamento de 18 de novembro da \researchq{3}. Essas execuções usaram várias
revisões da biblioteca com \texttt{gpt-5.6-terra} no nível padrão, antes da
configuração congelada (Seção~\ref{subsec:frozen}) e de várias regras de roteamento
(Seção~\ref{sec:pipeline}): um estudo de caso datado de uma versão anterior, não uma
medição da atual. \researchq{3} relata pull requests e issues
em 18 de novembro de 2026; \researchq{6}, a comparação da
Seção~\ref{subsec:protocol}; \researchq{7}, as reexecuções do RVSec sobre os pull
requests publicados, de 14 a 19 de novembro de 2026.

\subsection{\researchq{1}: Análise dos Achados e Possibilidade de Correção}
\label{subsec:rq1}

A adjudicação dos 34 achados registrou 17 decisões de propor uma correção
revisável (\texttt{MULTI\_FILE\_FIX\_PROPOSED\_FOR\_REVIEW}), 13 de que não
havia correção segura (\texttt{NO\_SAFE\_FIX}) e quatro de que um contrato
externo de digest impedia uma correção apenas no cliente. A Tabela~\ref{tab:adjudication} mostra
onde cada achado terminou, como registros locais: ali, uma issue é uma proposta mantida
para revisão; a publicação é uma etapa separada, fora das execuções do retrato.

\begin{table}[H]
\centering
\caption{Resultado dos 34 achados no retrato de outubro de 2026 (registros
locais).}
\label{tab:adjudication}
\footnotesize
\begin{tabular}{p{6.1cm}r}
\toprule
Resultado registrado & $n$ \\
\midrule
Correção gerada & 11 \\
\quad com commit na ponta do branch e build bem-sucedido\textsuperscript{a} & 7 \\
\quad aplicada; build falhou, como na base não modificada & 1 \\
\quad não aplicável ao branch base atual & 3 \\
Proposta local de issue para os mantenedores & 16 \\
Falso positivo reconhecido & 3 \\
Evidência insuficiente & 3 \\
Violação confirmada sem correção segura & 1 \\
\midrule
Total & 34 \\
\bottomrule
\end{tabular}
\\[2pt]\raggedright\scriptsize\textsuperscript{a}Seis até gerar um APK; no Feeder,
apenas o módulo da aplicação foi compilado.
\end{table}

As três correções não aplicáveis encontraram um branch base que já havia mudado,
em dois casos por pull requests mergeados deste próprio trabalho, e por isso
nenhuma produziu um pull request no retrato. A base do Nextcloud Cookbook continha
o \#204 (mergeado em 9 de agosto de 2026). O mantenedor do Giggity havia feito uma
mudança equivalente de forma independente e fechado o \#543. No DSub2000, o método
relatado havia sido alterado pelo \#160 (mergeado em 31 de agosto de 2026), que
tratou outra mensagem da mesma regra, mas a base ainda inicializava o TLS com um
trust manager que aceita tudo; uma verificação rejeitou a adaptação porque ela
removia o que a opção de conexão insegura do usuário ativava, e a etapa de
porte se absteve. Uma execução posterior abriu o DSub2000 \#178 para esse
trust manager em 22 de outubro; ele foi retirado no mesmo dia, depois que o
mantenedor respondeu que ignorar a verificação de certificados é o propósito
dessa opção, e não está entre os 49 pull requests.

Seis decisões de proposta terminaram como propostas de issue porque a geração foi
bloqueada: os rascunhos precisavam de credenciais de produção que o repositório não tem (MTG
Familiar), alteravam um algoritmo exigido por protocolo (dois no mOTP), alteravam
nomes de bancos de dados persistidos (RedReader), sobrescreviam um tamanho de chave RSA escolhido pelo chamador
(Droid SCEP) ou chamavam métodos verificadores inexistentes (BitBanana, cuja
migração chegou a um pull request por meio de seu achado irmão): uma decisão positiva
de planejamento é mais fraca que uma correção validada. As outras dez vieram de
decisões de ausência de correção segura ou de contrato externo; a biblioteca atual abre uma
issue para esses casos apenas com impacto de segurança confirmado
(Seção~\ref{sec:pipeline}), que, pelos registros do retrato, faltaria a nove deles.

Quinze dos 34 achados também têm um pull request publicado (16 pull requests;
Tabela~\ref{tab:sample}); em 18 de novembro de 2026, os pull requests de cinco
achados estavam incorporados, os de cinco fechados e os de cinco abertos,
contando o DSub2000 como incorporado pelo \#160. Nove desses pull requests não
são eventos posteriores: foram abertos a partir de análises anteriores dos
mesmos achados, antes do retrato (de 8 de agosto a 12 de outubro de 2026), e
seis já estavam mergeados ou fechados quando ele foi executado. Execuções
posteriores mudaram alguns resultados. O achado \texttt{CipherSpec} do FreeOTP,
sem correção segura no retrato, foi analisado novamente em sete execuções
posteriores, e o FreeOTP mergeou um teste (\#521) e depois a correção de
produção (\#524). O teste veio da primeira dessas execuções: sua
adjudicação deixou a causa em aberto, e uma regra determinística, não usada por
nenhuma execução do retrato, apoiou a causa na observação em tempo de execução
do RVSec e elevou a resposta à decisão apenas de revisão
\texttt{MULTI\_FILE\_FIX\_}\allowbreak\texttt{PROPOSED\_FOR\_REVIEW}. A correção
de produção veio de uma execução posterior, cuja adjudicação confirmou a causa. A proposta de issue do Droid SCEP tornou-se o pull request \#2. Seis issues
tratam desses achados. As do Critical Maps (\#379, depois fechada por nós, e
\#380, ainda aberta) e do OpenBible (\#190) foram abertas em agosto, antes do
retrato. MTG Familiar \#669 e DSub2000 \#177, esta sobre o trust manager do
DSub2000, vieram em 20 de outubro de uma execução de issues separada, fora do
experimento consolidado, e RedReader \#1394 veio em 21 de outubro.
Para o WiFiAnalyzer, o retrato recusou uma correção e não registrou impacto de segurança,
o motivo do mantenedor para fechar o \#618, mas ainda assim propôs uma issue, o que
o roteamento atual não faria; para o \texttt{sha1} do RedReader, propôs uma correção
cuja premissa o mantenedor havia rejeitado ao fechar o \#1373.

\textbf{Resposta à \researchq{1}.} No retrato de outubro, 17 dos 34 achados
tiveram uma decisão de propor correção e 11 chegaram a uma correção gerada: sete
tiveram commit na ponta do branch e build bem-sucedido, uma foi aplicada mas
falhou no build, como também o commit não modificado, e três não eram aplicáveis
porque a base havia mudado, duas delas por pull requests anteriores deste
trabalho. Dezesseis terminaram como propostas locais de issue
(Tabela~\ref{tab:adjudication}). Essas são classificações de uma versão anterior da biblioteca;
execuções posteriores e respostas de mantenedores mudaram várias delas.

\subsection{\researchq{2}: Aplicação e Build}
\label{subsec:rq2}

As 17 decisões de proposta produziram 11 correções geradas
(Tabela~\ref{tab:adjudication}): oito aplicadas na ponta do branch como prévias
locais de pull request e três não aplicáveis a ela. Sete prévias tiveram
commit e build bem-sucedido, seis até gerar um APK (no Feeder, apenas o código-fonte do módulo da aplicação
foi compilado). A oitava, MetadataRemover, falha no processamento de anotações
do Hilt com o mesmo erro no commit não modificado.

Uma auditoria das prévias encontrou problemas que nenhum build revela: o BetterUntis renomeia
arquivos de cache persistidos sem declarar essa decisão de compatibilidade, o PassNotes
mantém um verificador de hostname que aceita tudo em seu branch de debug, e o
\texttt{trustAllCerts} do Feeder valida certificados após o patch, embora mantenha
seu nome. No BitBanana, um candidato inicial compilava omitindo leitores de credenciais
afetados pelo novo verificador armazenado; a entrega de 20 de outubro (um
auxiliar e seis telas de credenciais) mantém a verificação legada, substitui um
verificador armazenado apenas após um desbloqueio bem-sucedido e gerou um APK. Nenhum
fluxo da aplicação foi testado; a \researchq{7} verifica apenas se a violação
relatada ainda aparece no método alterado pelo patch em tempo de execução.

\textbf{Resposta à \researchq{2}.} Das 11 correções geradas, oito foram aplicadas
na ponta do branch e sete delas tiveram build bem-sucedido, seis até gerar um APK;
a oitava também falha no commit não modificado, e três encontraram uma base que
já havia mudado. O build
não revelou os problemas que a auditoria encontrou em três prévias.

\subsection{\researchq{3}: Respostas dos Mantenedores e Motivos Relatados}
\label{subsec:rq3}

A Tabela~\ref{tab:pr-state} apresenta o estado dos 49 pull requests em 18 de
novembro de 2026; 67\% receberam resposta de alguém que não o autor. Os
incorporados incluem o BombusMod \#316, aplicado manualmente, e o FreeOTP \#521,
que apenas adiciona um teste (a mudança de produção, \#524, também foi mergeada).
Pela origem da execução de cada pull request (Seção~\ref{subsec:selection}), 8
vêm da campanha do grupo do RVSec (4 incorporados, 1 fechado, 3 abertos), 15 de
nossas análises de relatórios do RVSec anteriores a ela (6, 4, 5) e 26 de
nossas próprias filas do RVSec (11, 5, 10); a divisão é descritiva, pois as
filas favoreceram repositórios que aceitam pull requests externos.

\begin{table}[H]
\centering
\caption{Estado dos 49 pull requests do RVSec em 18 de novembro de 2026.}
\label{tab:pr-state}
\begin{tabular}{lr}
\toprule
Estado & PRs \\
\midrule
Incorporados (20 mergeados, 1 manualmente) & 21 \\
Fechados após revisão                  & 10 \\
Abertos (16 sem resposta)            & 18 \\
\midrule
Total                                & 49 \\
Receberam resposta de mantenedor       & 33 \\
\bottomrule
\end{tabular}
\end{table}

Das dez issues abertas para achados do RVSec (três para achados da campanha
do grupo, cinco de nossas análises anteriores e duas de nossas filas), duas,
abertas em setembro junto com
os pull requests do BombusMod e do wwwjdic, perguntavam se os mantenedores aceitariam
um pull request; nenhuma recebeu resposta, e nós fechamos ambas. As outras
oito descreviam um achado sem correção segura e perguntavam como tratá-lo; quatro
receberam resposta e duas continuam abertas. Três dessas respostas forneceram o que
a regra não tinha como saber (um fallback mantido por usabilidade, servidores
autoassinados dos usuários, um algoritmo de checksum definido por uma fonte remota), e o mantenedor do DNSToggle
corrigiu o achado em uma release.

\begin{table}[H]
\centering
\caption{Motivos apresentados para os 10 pull requests fechados sem
incorporação.}
\label{tab:closures}
\footnotesize
\begin{tabular}{p{5.8cm}r}
\toprule
Motivo declarado & PRs \\
\midrule
Digest usado fora de um papel de segurança (checksum, nome de cache ou de arquivo,
identificador) & 6 \\
Premissa do detector não se sustenta (o provedor já gera um IV aleatório) & 1 \\
A mudança apenas satisfaz a regra; o código existente é equivalente & 1 \\
Mudança equivalente já feita no repositório original & 1 \\
Projeto não aceita pull requests & 1 \\
\midrule
Total & 10 \\
\bottomrule
\end{tabular}
\end{table}

As respostas dos mantenedores acrescentam conhecimento da aplicação que falta a uma regra genérica
(Tabela~\ref{tab:closures}). Os seis fechamentos relativos a digest (WiFiAnalyzer, rssreader,
whoBIRD, BetterUntis, RedReader, InterTune) contestaram a substituição de MD5 ou SHA-1
usados como checksum, nome de arquivo de cache ou identificador. No entanto, o mesmo tipo de
mudança foi mergeado em outros projetos: no AMARadio, em nomes de arquivos de cache (\#122);
no ShortMessageSecure, em identificadores estáveis de lista, mudança que seu mantenedor considerou
criptograficamente irrelevante, mas uma melhoria razoável (\#89); e no NotallyX,
na verificação de hash de exportação (\#1167). Se tal mudança é bem-vinda é uma
decisão do mantenedor, não uma propriedade da regra; por isso, o {\system} ainda propõe
uma mudança correta para uma violação sem impacto de segurança e registra a
resposta.

Outras respostas mudaram o patch. No Crazyflie, a revisão redirecionou um patch de um
auxiliar de TLS para a remoção do auxiliar obsoleto; no ShortMessageSecure \#83,
a revisão aceita manteve \texttt{HmacSHA1} e tornou explícito o tamanho da chave.
O revisor do Session observou que, no Android, ``SSL'' e ``TLS'' selecionam a mesma
implementação. Entre os pull requests abertos, o mantenedor do Feeder objetou
que restringir protocolos em uma opção que os usuários ativam quebraria alguns sites,
e o mantenedor do prim-ftpd concordou em remover o RC4 após testes. As respostas
chegam a análises posteriores da mesma especificação (Seção~\ref{sec:pipeline}),
e duas tornaram-se verificações da biblioteca: um IV gerado pelo provedor é tratado como aleatório
(Amethyst \#4642), e remover protocolos de uma opção de compatibilidade ativada pelo usuário
conta como regressão (Feeder \#1210).

Três pull requests, de nossas análises de relatórios do RVSec anteriores à
campanha do grupo, que ainda incluíam regras genéricas de API, tratam de especificações não criptográficas (fechamento de streams no BombusMod
e no wwwjdic, um mapa sincronizado no ntfy); isso não estabelece cobertura de
famílias arbitrárias de APIs.

\textbf{Resposta à \researchq{3}.} Em 18 de novembro de 2026, 21 dos 49 pull
requests haviam sido incorporados, 10 fechados e 18 continuavam abertos
(Tabela~\ref{tab:pr-state}). A maioria dos fechamentos contestou a relevância de segurança de um
digest usado fora de um papel de segurança, enquanto mudanças comparáveis de digest foram
mergeadas em três outros projetos. Issues que relatavam um achado sem patch
obtiveram resposta em quatro de oito casos. A incorporação e o fechamento devem ser lidos
à luz desses motivos, e duas respostas que mudaram a premissa de um patch tornaram-se
verificações da biblioteca.

\subsection{\researchq{4}: Falsos Positivos e Violações sem Impacto de Segurança}

O retrato reconheceu três falsos positivos, cada um em uma
especificação diferente: Aegis (\texttt{GCMParameterSpecSpec}), mOTP
(\texttt{SecretKeySpecSpec}) e FreeOTP (\texttt{AlgorithmParametersSpec}),
um julgamento da análise, e não uma medida da precisão do detector. Oito
achados eram violações reais de regra sem impacto de segurança: os três
achados de \texttt{KeyStoreSpec} (Network Survey, Traffic Light, Tokn), que tratam
da sequência de chamadas que lê uma chave, e não de sua proteção; três dos 12
achados de \texttt{MessageDigestSpec} (dois no WiFiAnalyzer, um no Giggity),
que usam um digest fraco como identificador; um achado de \texttt{SSLContextSpec}
do Nextcloud Cookbook; e o achado de
\texttt{KeyPairGeneratorSpec} do Droid SCEP. Outros quatro achados de \texttt{MessageDigestSpec}
estão presos a um contrato externo que a aplicação não pode alterar sozinha
(dois no Password Monitor, um no OpenBible e um no Critical Maps); seu
impacto não foi estabelecido em três e foi confirmado em um.

Sete das oito violações sem impacto têm um pull request: Traffic
Light \#238, Tokn \#138 e Nextcloud Cookbook \#204 foram incorporados;
Network Survey \#154, WiFiAnalyzer \#618 e Giggity \#543 foram fechados; Droid
SCEP \#2 está aberto. O Network Survey \#154 foi fechado porque \texttt{getEntry}
e \texttt{getKey} são equivalentes para aquela chave. No Amethyst \#4642, o
mantenedor mostrou que o provedor já gera um IV aleatório quando
\texttt{init} recebe apenas a chave, um valor que o RVSec não consegue rastrear (o patch
alterava apenas a implementação desktop). A Tabela~\ref{tab:pr-spec} detalha todos
os pull requests por especificação.

\begin{table}[H]
\centering
\caption{Estado dos 49 pull requests do RVSec em 18 de novembro de 2026, por
especificação do achado.}
\label{tab:pr-spec}
\footnotesize
\setlength{\tabcolsep}{4pt}
\begin{tabular}{lrrrr}
\toprule
Especificação & Incorp. & Fechados & Abertos & Total \\
\midrule
\texttt{MessageDigestSpec} & 3 & 8 & 1 & 12 \\
\texttt{SSLContextSpec}    & 4 & 0 & 7 & 11 \\
\texttt{KeyStoreSpec}      & 5 & 1 & 1 & 7 \\
\texttt{CipherSpec}        & 4 & 0 & 2 & 6 \\
Outras dez ou sem registro\textsuperscript{a} & 5 & 1 & 7 & 13 \\
\midrule
Total & 21 & 10 & 18 & 49 \\
\bottomrule
\end{tabular}
\\[2pt]\raggedright\scriptsize\textsuperscript{a}O HTTP Shortcuts \#576 não tem
especificação registrada no levantamento; a recuperada para a \researchq{7} é a \texttt{SSLContextSpec}.
\end{table}

\textbf{Resposta à \researchq{4}.} O retrato reconheceu três falsos
positivos, em três especificações, e oito violações sem impacto de
segurança entre os 34 achados, seis delas em \texttt{KeyStoreSpec} e
\texttt{MessageDigestSpec}; entre os 49 pull requests,
\texttt{MessageDigestSpec} também responde por oito dos dez fechamentos, quatro
deles de achados do estudo de caso. Para digests, a
questão decisiva é para que o valor é usado, o que uma regra sobre o
nome do algoritmo não consegue responder, e os mantenedores a responderam de forma diferente em
projetos diferentes.

\subsection{\researchq{5}: Uso de Modelos e Custo}

Os rastros de prompts das execuções do retrato (19 e 20 de outubro de 2026;
11 execuções) registram 569 chamadas ao modelo enviadas ao provedor para os 34 achados,
com custo de US\$~35,71 a preços de tabela, uma mediana de US\$~0,77 por
achado; requisições
respondidas pelo cache de respostas não têm custo e não são contadas. Para as
chamadas contadas, esses valores são um limite superior: todas as chamadas foram atendidas no nível padrão,
antes do nível Flex, e a entrada em cache do \texttt{gpt-5.6-terra} é cobrada pela
tarifa integral de entrada. O estudo de caso não foi reexecutado com o modelo
atual: nas oito unidades da checagem de modelo (Seção~\ref{subsec:frozen}), o
\texttt{gpt-6.1-terra} custou US\$~2,43, contra US\$~6,31 do
\texttt{gpt-5.6-terra} a preços de tabela, para as mesmas 83 chamadas, e o nível
Flex reduz esse preço à metade, de modo que os valores do retrato superestimam o
que a configuração atual gasta. Eles também somam todas as execuções e
reexecuções de um achado dentro do retrato, enquanto a comparação relata uma
única execução de uma unidade, de modo que os dois valores não são diretamente
comparáveis. Execuções anteriores sobre os mesmos achados, de 30 de agosto de
2026 até o retrato (946 chamadas, US\$~75,50 a preços de tabela), e
reexecuções posteriores não estão incluídas. As 11 correções geradas custaram
US\$~13,19 (37\%), dos quais US\$~11,00 (31\%) foram das sete com commit na
ponta do branch, e as 16 propostas de issue US\$~17,98 (50\%). Achados com
evidência insuficiente foram os mais baratos (média de US\$~0,42), e falsos
positivos custaram em média US\$~0,82, perto da mediana. O custo se concentra
em achados reexecutados: os dois achados do BitBanana, cuja migração do
verificador armazenado foi revisada repetidamente, respondem por US\$~10,43. O custo da configuração
congelada é relatado junto com a comparação (\researchq{6}).

\textbf{Resposta à \researchq{5}.} No retrato de outubro, com o
\texttt{gpt-5.6-terra} a preços de tabela (um limite superior), analisar um achado
custou uma mediana de US\$~0,77; as 16 propostas locais de issue ficaram com
metade do gasto e as 11 correções geradas com 37\%. Na
configuração congelada usada na comparação (\texttt{gpt-6.1-terra}, nível
Flex), uma execução de uma unidade de análise custou em média US\$~0,10, ou US\$~0,19 a
preço de tabela; este é um valor por execução, não comparável à mediana por
achado acima.

\subsection{\researchq{6}: Comparação com Agentes de Programação e Variação entre Execuções}
\label{subsec:rq6}

A comparação da Seção~\ref{subsec:protocol} foi executada sobre 20 unidades de análise (24 achados): nove
unidades sem pull request nosso, três cujo pull request um mantenedor
havia fechado e oito cujo pull request está aberto ou foi incorporado.
A Tabela~\ref{tab:optiona} apresenta as entregas de cada braço no commit analisado, o
destino de cada pull request na ponta do branch e as duas leituras dos juízes.

\begin{table*}[t]
\centering
\caption{\system{} e cinco braços de agentes de programação (letras como na
Seção~\ref{subsec:protocol}) em 20 unidades de análise (24 achados do RVSec), quatro execuções por unidade. Primeiro
bloco: agentes que receberam o achado e o repositório, julgados no commit analisado. Segundo
bloco: as mesmas configurações recebendo o processo do \system{} como instruções, com o
porte para a ponta do branch. As execuções aceitáveis são calculadas como média sobre
as unidades; uma issue ou uma nota de que nada deveria ser enviado conta como aceitável. A
leitura completa julga o cartão inteiro; a leitura do patch, apenas o patch. Custo: preço de
tabela para todos os braços; o \system{} foi cobrado em US\$~0,10 por execução no nível Flex.}
\label{tab:optiona}
\footnotesize
\setlength{\tabcolsep}{4pt}
\begin{tabular}{p{4.2cm}r|rrr|rr}
\toprule
 & & \multicolumn{3}{c|}{Primeiro bloco} & \multicolumn{2}{c}{Segundo bloco} \\
 & \system{} & Codex Terra & Claude Sonnet & Codex Sol & Codex Terra & Claude Sonnet \\
 & C & T & S & O & TP & SP2 \\
\midrule
Execuções & 80 & 80 & 80 & 80 & 80 & 80 \\
PR / issue / nada / falha & 31/10/39/0 & 23/8/49/0 & 16/13/51/0 & 15/2/63/0 & 38/16/24/2 & 46/5/27/2 \\
Unidades com PR (de 20) & 9 & 8 & 6 & 5 & 10 & 15 \\
\multicolumn{7}{l}{\emph{Pull requests na ponta do branch}} \\
\quad Portados e com build & 20 & --- & --- & --- & 24 & 34 \\
\quad Já corrigidos ou código removido & 11 & --- & --- & --- & 2 & 1 \\
\quad Abstenção ou falha & 0 & --- & --- & --- & 12 & 11 \\
\multicolumn{7}{l}{\emph{Juízes}} \\
\quad Aceitáveis, leitura completa & 0,96 & 0,94 & 0,97 & 0,90 & 0,85 & 0,80 \\
\quad Aceitáveis, leitura do patch & 0,93 & 0,99 & 0,97 & 1,00 & 0,85 & 0,70 \\
\quad PRs propostos & 31 & 23 & 16 & 15 & 38 & 46 \\
\quad PRs aceitáveis, completa / patch & 28 / 25 & 18 / 22 & 14 / 14 & 7 / 15 & 28 / 28 & 32 / 24 \\
\quad Unidades com PR aceitável, completa / patch & 8 / 7 & 6 / 7 & 5 / 5 & 3 / 5 & 8 / 8 & 11 / 8 \\
Custo por execução, tabela (US\$) & 0,19 & 0,04 & 0,12 & 0,13 & 0,21 & 0,37 \\
\bottomrule
\end{tabular}
\end{table*}

Os 11 pull requests do \system{} cuja ponta já continha a correção ou havia perdido
o código (AMARadio, ANOTHERpass, NotallyX) não passaram por build e são julgados
apenas pelo patch;
sobre os 29 que o \system{} publicaria (ele reteria os outros dois, cujo patch não
tratava um achado reportado), as leituras são 0,95 e 0,91. Dos 84
do segundo bloco, 57 foram aplicados sem conflito na ponta, 5 foram portados após um
conflito e 22 terminaram como já corrigidos ou em abstenção; os quatro do Codex Terra portados
após um conflito não tiveram build bem-sucedido (o sandbox do agente não conseguia gravar o keystore
de debug), e seus patches já haviam sido julgados inaceitáveis.

\paragraph{(a) Um agente que recebe o achado e o repositório.} O Claude Sonnet
pontuou acima do \system{} nas duas leituras; o Codex Terra e o Codex Sol pontuaram
abaixo dele na leitura completa e acima na leitura do patch. Os
três propuseram menos pull requests (23, 16 e 15 contra 31), com menos
julgados aceitáveis na leitura completa (18, 14 e 7 contra 28). Eles não portaram
seus patches para a ponta, onde um pull request do \system{} ainda pode falhar, nem
receberam a política do \system{} de que uma alteração correta apenas de regra vale a pena ser proposta.

\paragraph{(b) Um agente instruído a seguir o mesmo processo.}
Recebendo o processo como instruções, os dois agentes propuseram mais pull requests
que o \system{} (38 e 46 contra 31) e tiveram tantos ou mais julgados aceitáveis
na leitura completa (28 e 32 contra 28), em tantas ou mais unidades (8 e 11
contra 8); na leitura do patch, o braço do Claude Sonnet teve menos (28 e 24
contra 25). Na leitura completa, também tiveram mais julgados inaceitáveis
(propostos menos aceitáveis): 10 do Codex Terra e 14 do Claude Sonnet, contra 3
do \system{}. Sua proporção de execuções aceitáveis foi menor (0,85 e 0,80 contra
0,96 na leitura completa; 0,85 e 0,70 contra 0,93 na leitura do
patch), e cada um teve duas execuções com falha.
O \system{} e os braços Codex Terra compartilham o modelo e, no segundo bloco,
as entradas e as regras; diferem na execução. O \system{} executa cada etapa como uma
chamada separada ao modelo, decide se a violação é real e se existe uma correção segura
antes de qualquer rascunho, oculta a extração da adjudicação e os invariantes do plano
de quem redige o rascunho, e impõe verificações e novas tentativas em código; o agente
segue as etapas em uma única conversa e verifica o próprio trabalho.
O Codex Terra com as instruções encontrou tantos pull requests aceitáveis quanto o
\system{}, mas teve mais julgados inaceitáveis; a diferença é compatível com a
execução em etapas e com uma etapa de decisão separada da edição, e não com as
instruções, que ambos tinham, embora 20 unidades não permitam estabelecê-lo.

O custo de tabela do \system{}
(US\$~0,19 por execução) inclui suas buscas na web e o porte, assim como o do segundo bloco
(US\$~0,21 e US\$~0,37) inclui seu porte.

\paragraph{Variação entre execuções.} O \system{} executou cada unidade quatro vezes sem o
cache de respostas. A classe de entrega (pull request, issue ou nenhuma ação) foi a
mesma nas quatro réplicas em 16 das 20 unidades, e o desfecho registrado, em
11. A classe mudou para o HMAC do mOTP e a derivação de chave do Salty (pull request
ou nenhuma ação) e para a derivação de chave do Winterhold e a configuração TLS do Read You (pull
request ou issue); em quatro unidades mudou apenas o motivo (por exemplo, nenhuma correção segura vs.\
contrato externo), e o NotallyX alternou entre uma correção que a ponta já continha e um
pull request com build bem-sucedido. O custo cobrado variou até 3,8 vezes entre as réplicas
(Read You, de US\$~0,074 a US\$~0,28, sendo a mais cara encerrada como issue).
A classe dos agentes foi a mesma nas quatro réplicas em 12 a 17 unidades
(Codex Terra 12, Claude Sonnet 13, Codex Sol 17, e 17 e 14 no segundo
bloco). Uma única execução amostra uma
decisão; diferenças de uma ou duas unidades entre braços estão dentro dessa
variação.

\paragraph{Juízes e cegamento.} Os juízes concordaram quanto à aceitabilidade com
$\kappa = 0{,}81$ (491 cartões, 15 deles da execução parcial do \system{} com o Claude
Sonnet 5.5) e quanto à nota do patch com um $\kappa$ ponderado quadraticamente
de $0{,}82$ (174 pull requests, 5 dessa execução) na leitura completa. Na sondagem do primeiro bloco
(16 de novembro de 2026), nos cartões completos o juiz principal apontou o \system{} em
todos os 8 cartões dele, mas também em 13 dos 24 cartões dos três braços de agentes e em 4
dos 8 cartões da execução parcial do \system{} com o Claude Sonnet 5.5, com confiança média
entre 0,28 e 0,35: um palpite padrão, não uma identificação. Nos cartões apenas com o patch, apontou o
\system{} em 1 de 10. Na sondagem do segundo bloco (10 cartões por braço), apontou o \system{} em 23 de
30 cartões completos, 9 deles do \system{}, e o Codex Terra em 21 de 30 cartões apenas com o patch,
8 deles do Codex Terra, com confiança média entre 0,38 e 0,53:
novamente um palpite padrão.

\textbf{Resposta à \researchq{6}.} (a) Agentes que receberam apenas o achado e o
repositório propuseram menos pull requests que o \system{} (23, 16 e 15 contra 31),
com menos julgados aceitáveis na leitura completa (18, 14 e 7 contra 28). Um escore
por execução que conta issues e desfechos sem ação como aceitáveis, e assim recompensa
a abstenção, os colocou perto do \system{} ou acima dele (0,90 a 0,97, contra
0,96). (b) Agentes instruídos a seguir o mesmo processo propuseram mais pull
requests (38 e 46) e tiveram tantos ou mais julgados aceitáveis na leitura completa
(28 e 32, contra 28), mas mais julgados inaceitáveis (10 e 14, contra 3 de 31 do
\system{}); o mesmo escore os colocou abaixo dele (0,85 e 0,80). Com o mesmo modelo
e as mesmas instruções, a diferença é compatível com executar o processo em
etapas separadas, com a decisão de agir tomada antes de qualquer rascunho. Ao
longo de suas quatro execuções por unidade, o \system{} entregou o mesmo tipo de
resultado em 16 de 20 unidades.

\subsection{\researchq{7}: RVSec nos APKs com patch}
\label{subsec:rq7}

Nas seis prévias do retrato que geraram um APK (três repetições cada), a violação
reportada deixou de aparecer no método alterado pelo patch em cinco: Tokn, Traffic Light,
Network Survey, BetterUntis e PassNotes (cujo APK original a exibiu em uma de
três; isso verificou o patch da própria execução, que substitui um trust manager
que aceita tudo, e não o keepassvault \#347 publicado, relatado abaixo). BitBanana foi inconclusivo: seu método só é executado com um bloqueio do app
configurado, alcançado em uma de três execuções do APK original e em nenhuma do APK com patch.
A Tabela~\ref{tab:dynval} reporta todos os 49 pull requests: os 45 executados conforme a
Seção~\ref{subsec:dynval}, o MetadataRemover \#264, cujo repositório não
compila no commit base do pull request (um plugin do Hilt mais antigo que as
bibliotecas do Hilt trazidas por uma atualização de dependência; o build do
patch feito pelo próprio pipeline falhou na mesma etapa, enquanto o commit analisado compilou), e três cuja especificação está fora do conjunto
do RVSec.

\begin{table}[H]
\centering
\caption{Reexecução do RVSec no APK original e no APK com patch dos 49 pull requests
publicados, por estado em 18 de novembro de 2026. ``Corrigido'': a violação reportada
deixou de aparecer no método alterado (incluindo a remoção da chamada
governada, um caminho que o patch desativa e uma mensagem reportada que sumiu, restando no método
apenas mensagens que o APK original já exibia). ``Não verif.'':
a violação reportada não apareceu em nenhum dos APKs e sua mensagem não é mais
emitida. ``Testes'': o pull request altera apenas testes. Dos 38 verificados (Corr., Não corr., Testes),
33 estão corrigidos (25 pela regra registrada, sem os rótulos acrescentados depois das execuções).}
\label{tab:dynval}
\footnotesize
\setlength{\tabcolsep}{2pt}
\begin{tabular}{lrrrrrrr}
\toprule
 & Corr. & Não corr. & Testes & Não verif. & Incon. & Não exec. & Total \\
\midrule
Incorporados & 14 & 2 & 1 & 3 & 0 & 1 & 21 \\
Fechados     & 8 & 2 & 0 & 0 & 0 & 0 & 10 \\
Abertos      & 11 & 0 & 0 & 3 & 1 & 3 & 18 \\
\midrule
Total        & 33 & 4 & 1 & 6 & 1 & 4 & 49 \\
\bottomrule
\end{tabular}
\end{table}

Pela origem da execução por trás de cada pull request
(Seção~\ref{subsec:selection}), os 26 das nossas filas do RVSec, que usaram os
exploradores da reexecução, ficaram em 22 corrigidos e 4 não corrigidos; os 8 da
campanha do grupo do RVSec, em 5 corrigidos, 1 só com testes, 1 inconclusivo e 1
não executado; e os 15 das nossas execuções anteriores sobre dados do RVSec
anteriores à campanha, em 6 corrigidos, 6 não verificáveis e 3 não executados.
Os quatro não corrigidos vêm todos das filas, e os seis não verificáveis, todos
das execuções anteriores. A divisão é descritiva: os grupos diferem em como e
quando o RVSec foi executado.

A incorporação e a reverificação em tempo de execução divergem nos dois sentidos. Alguns pull requests
incorporados mantiveram a violação reportada: o ShortMessageSecure \#83 manteve \texttt{HmacSHA1} a pedido do mantenedor, de modo que
a regra do gerador de chaves continua violada, e no ANOTHERpass \#103 a sequência de chamadas
do KeyStore continua sendo reportada. Em cinco pull requests, dois incorporados (Session \#2241,
Nextcloud Cookbook \#204) e três abertos (Fread \#168, zSMTH \#66, Grocy
\#1012), a mensagem reportada desapareceu; outras mensagens da mesma especificação, já presentes no
APK original e não modificadas pelo patch, permaneceram na mesma chamada.
Inversamente, a violação desapareceu em seis pull requests cujos mantenedores
objetaram à mudança (whoBIRD, rssreader, RedReader, BetterUntis, Network
Survey, WiFiAnalyzer): satisfazer o monitor não tornou essas mudanças
desejadas. Ela também desapareceu no Giggity \#543, fechado porque o mantenedor
já havia feito uma mudança equivalente, e no Summit \#54, fechado porque o
projeto não aceita pull requests. No Amethyst \#4642, fechado pelos motivos dados na \researchq{4},
a mensagem permaneceu, pois o patch alterou apenas a implementação desktop,
que o APK não inclui; no InterTune \#52 o digest SHA-1 passou para outra função do
mesmo arquivo, onde o RVSec continuou a reportá-lo. No Read You \#1334 o método
alterado deixou de ser executado porque o patch desativou, por padrão, a opção que
levava a ele.
Com o DroidBot, o explorador do relatório original, a violação do FreeOTP
apareceu no APK original em todas as três repetições de ambos os pull requests;
ela desapareceu após o \#524, a mudança de produção, e permaneceu após o \#521,
que apenas adiciona um teste.

Seis relatórios não puderam ser verificados com o RVSec atual. No News \#1716 e no
keepassvault (PassNotes) \#347 o APK original executou o método reportado sem nenhuma
violação; no DSub2000 \#160, Feeder \#1210, etesync \#301 e HTTP Shortcuts
\#576 a violação tratada pelo pull request não apareceu em nenhum dos APKs, e
o que restou na chamada (um key manager ou fonte de aleatoriedade nulos, um
trust manager definido pela aplicação) já estava no APK original e fora
do pull request. Os seis relatórios trazem mensagens de protocolo ou de sequência de chamadas que
nenhuma execução da validação emitiu.

Uma segunda leitura adversarial de 21 diffs rotulados como corrigidos e a
leitura de um único pesquisador de outros seis (três cuja mensagem reportada
desapareceu e três acrescentados a partir do formato de mensagem anterior)
tentaram mostrar que o patch apenas realocava ou encapsulava a chamada. Os
outros seis casos corrigidos (Your Local Weather \#209, Read You \#1334,
Session \#2241, Nextcloud Cookbook \#204 e Summit \#54, rotulados depois dessas
leituras, e FreeOTP \#524, lido abaixo) não foram classificados. Dos 27, catorze alteram o
valor ou material governado (um algoritmo de digest, um tamanho de chave, um protocolo, os certificados
confiáveis, a origem de um IV ou de um salt, ou a remoção do RC4), nove alteram a sequência de chamadas ou a API para a
forma exigida com o mesmo material, quatro removem a chamada governada e nenhum
apenas a realoca; nenhuma execução com patch reportou a mesma especificação em um novo
local. Em oito, o comportamento criptográfico não muda: esses patches levam o código à
forma que a regra exige, como faz a correção de uma violação sem impacto
de segurança. Nos casos de KeyStore o relatório desaparece porque o monitor
passa a ver de onde a chave veio, não porque vê menos. Um patch
corrigido, o FreeOTP \#524, mantém o mesmo material de chave: ele reconstrói a
chave derivada da senha como uma chave AES a partir dos próprios bytes e acrescenta uma verificação de que
a chave é AES, de modo que o monitor deixa de reportar a origem da chave enquanto a
cifração permanece inalterada. Esse é o padrão que a revisão de comportamento deveria
rejeitar (Seção~\ref{sec:pipeline}), e ele chegou a um pull request mergeado.
Um comentário no monitor \texttt{CipherSpec} do RVSec, junto à verificação que
emite esse relatório, aponta essa reconstrução como o caminho que as regras CrySL
admitem para uma chave derivada de senha usada com AES.

\textbf{Resposta à \researchq{7}.} Ao reexecutar o RVSec nos pull requests
publicados, a violação reportada deixou de aparecer no método alterado em
33 dos 38 que ele pôde verificar (25 pela regra registrada); ainda apareceu
em quatro, e um pull request
alterou apenas testes. Seis relatórios trazem mensagens que o RVSec atual não
emite mais e não puderam ser verificados, e em um caso o método alterado nunca foi
alcançado. Um patch corrigido apenas mudou a forma como o monitor via a mesma chave. Um build, o desaparecimento da violação reportada e a
aceitação por um mantenedor são propriedades distintas, e os casos mostram cada uma
sem as demais.

\section{Discussão}
\label{sec:discussion}

\subsection{Principais Achados}

Quatro observações resumem o estudo. Primeiro, os mantenedores julgaram os achados em
contexto: dos 49 pull requests, 21 foram incorporados e 10 foram fechados após
revisão (Seção~\ref{sec:results}). Suas respostas levantaram o comportamento
intencional, a compatibilidade com valores armazenados e a relevância da regra
para eles, conhecimento ausente da regra e da chamada reportada. Uma mudança revisável fez
esse conhecimento emergir com mais frequência do que uma issue: 33 dos 49 pull
requests e quatro das dez issues receberam resposta (Seção~\ref{subsec:rq3}; uma
comparação descritiva).

Segundo, uma violação real não exige necessariamente um reparo de segurança. Entre os 34
achados consolidados em outubro, os relatórios sem impacto de segurança concentraram-se
em digests de mensagem e key stores, onde o uso do valor é decisivo
(\researchq{4}); seis dos dez pull requests fechados substituíam MD5 ou SHA-1 que
o mantenedor via como checksum, nome de cache ou outro valor sem
papel de segurança.

Terceiro, um patch que compila não necessariamente remove a violação.
Ao reexecutar o RVSec nos APKs original e com patch, a violação reportada deixou
de aparecer no método alterado em 33 dos 38 casos verificáveis
(25 pela regra registrada), ainda apareceu em quatro, e um pull request apenas adicionou um teste
(Seção~\ref{subsec:rq7}). Dos quatro, um manteve uma escolha que o mantenedor pediu
para preservar, um deixou uma sequência de chamadas que o monitor ainda rejeita, um alterou
apenas uma implementação desktop e um moveu o digest para outra função do
mesmo arquivo.

Quarto, agentes que receberam apenas o achado e o repositório propuseram menos pull
requests que o \system{} e tiveram menos julgados aceitáveis; só um escore por
execução que conta a abstenção como aceitável os colocou próximos ou acima dele (0,90 a
0,97, contra 0,96, na leitura completa; \researchq{6}). Agentes que receberam o processo como
instruções tiveram tantos ou mais pull requests julgados aceitáveis quanto o
\system{} na leitura completa, mas de três a cinco vezes mais julgados inaceitáveis; a diferença é compatível com a forma como o
processo é executado (etapas separadas, cada uma verificada por código, com uma decisão
antes de qualquer edição) e não com as instruções em si, embora 20 unidades não
permitam estabelecê-lo.

\subsection{Como uma Contribuição É Lida}
\label{subsec:form}

A forma das primeiras contribuições suscitou tantos comentários quanto seu conteúdo. Os pull
requests até meados de setembro seguiam um template de relatório do pipeline: um título como
``Fix \texttt{<init>}'', de três a onze títulos de seção, checklists e até
cerca de 950 palavras. Em outubro o pipeline passou a escrever para o mantenedor:
um resumo curto sob dois títulos e, a partir de 20 de outubro, alguns parágrafos
simples sobre o que o código fazia, o que mudou e que o build passa, com
títulos apenas quando o template do repositório pede (Tabela~\ref{tab:pr-format}).
O exemplo completo da Seção~\ref{sec:pipeline} (\#238) é um desses pull requests
em texto simples.

\begin{table}[H]
\centering
\caption{Formato dos 49 pull requests do RVSec e seu estado em 18 de novembro de
2026. Palavras e títulos são medianas do corpo atual, arredondadas.}
\label{tab:pr-format}
\footnotesize
\setlength{\tabcolsep}{1.5pt}
\begin{tabular}{lrrr}
\toprule
 & \begin{tabular}[b]{@{}r@{}}Com template\\8/8--16/9\end{tabular}
 & \begin{tabular}[b]{@{}r@{}}Seções curtas\\12/10--20/10\end{tabular}
 & \begin{tabular}[b]{@{}r@{}}Texto simples\\20/10--7/11\end{tabular} \\
\midrule
Pull requests & 15 & 12 & 22 \\
Palavras & 227 & 207 & 157 \\
Títulos de seção & 6 & 2 & 0 \\
Títulos ``Fix \textit{método}'' & 12 & 0 & 0 \\
Incorporados & 6 & 4 & 11 \\
Fechados após revisão & 4 & 4 & 2 \\
Abertos & 5 & 4 & 9 \\
Receberam resposta & 11 & 8 & 14 \\
Objeção ao texto ou à origem & 3 & 0 & 0 \\
\bottomrule
\end{tabular}
\end{table}

O mantenedor do
BetterUntis recusou contribuições que considerava geradas por IA, o mantenedor do Feeder
se recusou a ler ``that wall of ai text,'' e o revisor do Crazyflie
pediu um título e uma descrição explicando o que foi corrigido e por quê.
Essas objeções se dirigem ao texto ou à sua origem, não à mudança, e todas vieram
do período com template; nenhum dos 34 pull requests posteriores
havia recebido alguma até 18 de novembro de 2026, e os seis deles fechados após revisão
foram fechados por razões técnicas ou por uma política de contribuição.

Um texto que parece saída de ferramenta---um título genérico, vocabulário interno,
longas seções de evidências---sugere que ninguém revisou a mudança, de modo que o
mantenedor questiona sua origem antes de seu conteúdo; um texto curto nos
termos do próprio projeto permite que a revisão comece pelo código. Nada é ocultado: 32 dos 49
textos nomeiam o RVSec como analisador, e os outros 17, abertos de 8 de agosto a
22 de outubro, nomeiam o explorador que conduziu a interface (14) ou nenhuma
ferramenta (3); o autor aprova cada submissão e o uso de IA é declarado onde um
template pede. A comparação tem fatores de confusão. Os
períodos diferem na origem dos achados: os 15 pull requests com template vieram
todos de nossas execuções de agosto e setembro sobre dados do RVSec anteriores à
campanha do grupo, enquanto 18 dos 22 em texto simples vieram de nossas próprias
filas, parte deles de uma fila que só admitia repositórios que haviam incorporado
um pull request externo nos doze meses anteriores. Os pull requests posteriores
também vieram de uma análise mais madura e de outros projetos, os mais recentes
tiveram menos tempo para revisão (parte de sua contagem de abertos), e os corpos
foram lidos após quaisquer alterações feitas na revisão (o do Crazyflie foi
reescrito conforme solicitado). A tabela não mostra que o formato aumenta a
aceitação; o que mudou foi o tipo de objeção (os fechamentos restantes dizem
respeito à mudança em si), e essa mudança pode refletir também a origem dos
achados e os repositórios que os receberam.

\subsection{Implicações}

Para pesquisadores, avaliações de reparo de segurança deveriam contar a abstenção como um
resultado, reportando-a ao lado da validade do patch (o \system{} registra falsos
positivos, violações sem impacto de segurança ou sem correção segura, contratos
externos, correções já presentes no upstream, issues e pull requests), e deveriam dar
a um pipeline e a um agente a mesma tarefa (os agentes do nosso primeiro bloco pararam no
commit analisado; o segundo bloco também faz o porte). A completude das evidências
merece uma métrica própria: a migração do verificador armazenado deixou leitores de fora mesmo
com o alvo localizado e um candidato compilável. Benchmarks também deveriam
manter e atribuir as falhas de infraestrutura (JDK, SDK, desatualização de plugins) em vez
de descartar projetos que não compilam.

Para praticantes, o fluxo de trabalho é um assistente de revisão, não um autor
autônomo de commits: a execução registra por que um alerta importa, quais suposições permanecem
e que validação falta, enquanto o texto para os mantenedores permanece curto.
A publicação fica fora do ciclo de análise: nenhuma credencial em prompts, nenhuma
resposta do modelo capaz de publicar, aprovação vinculada a um identificador de achado e a uma
validação bem-sucedida, e política e testes dedicados para mudanças na aceitação de
certificados, na verificação de hostname ou em formatos criptográficos.

O custo é uma referência, não o fator decisivo: os custos por execução foram de US\$~0,04 a
US\$~0,37 a preço de tabela, entre o \system{} e os agentes
(Tabela~\ref{tab:optiona}), e a escolha depende da qualidade do que é enviado
(Seção~\ref{subsec:rq6}).

\paragraph{Pipeline ou agente?} A comparação sugere que o que importa é
executar o processo em etapas separadas, com a decisão de agir tomada antes de
qualquer edição (Seção~\ref{subsec:rq6}). Além disso, o \system{} mapeia um
relatório de tempo de execução para seu código-fonte, une
achados no mesmo código, mantém as respostas anteriores dos mantenedores e submete cada
patch a verificações determinísticas, a um build, ao porte e a uma revisão de comportamento
antes que uma pessoa decida. Seu modelo devolve edições como dados e nunca executa
comandos, mas os builds executam os próprios scripts Gradle de cada repositório na
estação de trabalho do autor, sem contêiner (Seção~\ref{subsec:ethics}). Os
agentes da comparação foram executados ali com ferramentas restritas e nenhum segredo além do
login da CLI; um agente com shell livre sobre centenas de repositórios desconhecidos precisa de um
contêiner descartável por execução que veja apenas o repositório, alcance apenas
registros de pacotes e a API do modelo, e não guarde segredos. As abordagens
são complementares: o \system{} prepara um achado e decide o que chega a um
mantenedor, enquanto um agente explora e aplica patches dentro de tal ambiente
(Trabalhos Futuros). Modelos locais via Ollama são suportados, sem pesquisa
na web, mas não foram avaliados.

\subsection{Lições Aprendidas}
\label{sec:lessons}

\textbf{Alcançar o código é o primeiro gargalo.} O RVSec reporta apenas o código
que a exploração executa. Em nossas próprias filas, o Monkey e o APE raramente
passaram de telas de login ou de configuração, de modo que a criptografia protegida
por uma conta ou por uma configuração de servidor fica sub-representada entre seus
achados. Isso também limita a verificação de um patch: nas reexecuções, apenas o
DroidBot, um dos exploradores do relatório original, alcançou o método reportado do FreeOTP, e o método
modificado pelo patch do BitBanana nunca foi alcançado (\researchq{7}). Os exploradores decidem quais achados existem e quais
correções podem ser confirmadas. O grupo do RVSec mediu o mesmo limite em 163
aplicações: na aplicação mediana com tais métodos, a exploração exercitou cerca de
1,4\% dos métodos que chamam diretamente uma API monitorada~\cite{costa2026effectiveness}.

\textbf{Localizar o código faz parte do reparo.} O RVSec nomeia um método de
bytecode, não uma função do código-fonte. Lambdas, corrotinas e nomes alterados
pelo compilador Kotlin tiveram de ser rastreados até a função que usa a API e,
em Kotlin Multiplatform, até a implementação Android: o Amethyst \#4642 alterou
apenas a implementação desktop, código que o APK nunca contém. Um alvo errado
produz um patch que compila e deixa intocado o código reportado; por isso, uma
correspondência ambígua fica sem resolução em vez de ser adivinhada.

\textbf{Uma violação não é um defeito até que se conheça o seu uso.} MD5 como chave
de cache, uma chave lida com \texttt{getEntry} em vez de \texttt{getKey} ou um contexto
SSL com um trust manager próprio para um servidor autoassinado violam uma regra sem
expor o usuário; os mantenedores rejeitaram alterações, ou responderam a issues, que
ignoravam isso, e o que decidiu cada caso foi
a forma como o valor era consumido, não a chamada de API.

\textbf{Verificações determinísticas são um seguro barato, se enxergarem o mesmo contexto.}
No retrato de outubro, verificações semânticas recusaram seis das 17 propostas: cinco
das seis gerações bloqueadas da \researchq{1} e a adaptação do DSub2000 no porte. A
auditoria considerou corretas quatro recusas: rascunhos que substituíam dois
algoritmos exigidos por protocolo no mOTP (MD5 para os códigos mOTP, HmacSHA1
para TOTP), alteravam os nomes de banco de dados armazenados
do RedReader e sobrescreviam o tamanho de chave escolhido pelo chamador no Droid
SCEP. As outras duas vieram de defeitos do pipeline: a verificação de exceções leu
apenas o fragmento editado do rascunho do MTG Familiar e não viu um \texttt{catch}
envolvente (um bloqueio indevido, embora aquele rascunho também dependesse
de credenciais de produção que o repositório não tem), e a verificação da opção não havia lido uma alteração que a
adjudicação autorizara para o DSub2000.
As verificações removem patches plausíveis porém prejudiciais sem custo de modelo, desde que cada uma
leia o contexto da etapa que verifica.

\textbf{Uma reexecução é lida restrição por restrição.} O RVSec monitora uma
especificação inteira, de modo que um patch pode remover a restrição reportada e
deixar outra, e um mantenedor pode manter o que uma regra rejeita; por isso, o
veredito de uma reexecução é lido por restrição, ao lado da decisão do mantenedor
(\researchq{7}).

\textbf{Alterações pequenas são revisadas; grandes são discutidas.} Alterações de uma
única chamada com poucas frases de explicação costumavam ser resolvidas sem
debate, enquanto renomeações, configurações ou migrações de dados armazenados suscitavam questões de
escopo (Crazyflie, ShortMessageSecure, RedReader, InterTune).

\textbf{O alvo é a ponta do branch, não o commit analisado.} Os achados vêm de
um APK publicado, enquanto os mantenedores revisam sobre o branch atual. O porte
para ele encontrou correções já presentes no upstream ou arquivos reportados
removidos, casos em que o desfecho correto era parar; um pull request foi
fechado porque uma mudança equivalente já havia sido feita no upstream
(Tabela~\ref{tab:closures}); e o MetadataRemover \#264 foi aberto embora não compile no seu
commit base, falha que o branch sem modificação também tem e que o texto do pull
request declara, enquanto o commit analisado compilava (\researchq{7}).

\textbf{O texto é lido antes da mudança.} As objeções ao texto ou à origem
de um pull request, e não à mudança, vieram apenas do período com template, que
também difere dos posteriores na origem de seus achados
(Seção~\ref{subsec:form}). As políticas do repositório fazem
parte do mesmo contexto: restrições a contribuições assistidas por IA ou regras
que admitem apenas issues foram mostradas ao autor como informação, não como
bloqueios automáticos, e a decisão ficou com uma pessoa.

\subsection{Considerações Éticas}
\label{subsec:ethics}

Os pull requests e issues foram enviados a mantenedores voluntários, cujo tempo
consumiram e cujas respostas públicas são dados deste trabalho. O autor aprovou cada
submissão após ler seu diff e seu texto. Um pull request só foi aberto para
um patch que se aplicava e passava no build, exceto o MetadataRemover \#264, cujo build
falhou na mesma etapa no branch sem modificação, como seu texto informa. Dos
49 textos de pull request, 32 nomeiam o RVSec como analisador e os outros 17 nomeiam o
explorador que conduziu a execução ou nenhuma ferramenta; o uso de IA é declarado quando um
template o pergunta, e a política de IA de um projeto era mostrada ao autor, que então
escrevia o texto. Antes do estudo, ao testar o pipeline com um
modelo local, abrimos 23 issues em 4 de agosto de 2026 com o mesmo texto curto,
perguntando se um pull request seria bem-vindo; elas ficam fora de todas as contagens. As reações dos mantenedores aos pull
requests com template nos levaram a encurtar os textos (Seção~\ref{subsec:form}). Todos os dados são registros públicos do GitHub; o artigo
nomeia projetos e pull requests, não mantenedores. Os builds de terceiros rodaram na
estação de trabalho do autor sem contêiner, ao lado das credenciais de
publicação; isolar cada build é trabalho futuro.

\subsection{Ameaças à Validade}
\label{subsec:threats}

\subsubsection{Validade interna}

\emph{Implementação em evolução.} O pipeline, os prompts e as verificações mudaram durante
o estudo, de modo que o retrato de outubro com 34 achados é um estudo de caso datado, oriundo de
várias versões do \system{} (resultado mais recente por achado), e não o desempenho
de uma única configuração. Os pull requests publicados também vieram de diferentes
versões, alguns alterados durante a revisão ou portados, e cada desfecho é atribuído à
revisão submetida. Apenas a comparação com agentes de
programação usou uma única configuração registrada (Seção~\ref{subsec:frozen}).

\emph{Não determinismo do modelo.} Na comparação, o \system{} analisou cada uma
das 20 unidades quatro vezes sem o cache de respostas, e sua classe de entrega
(pull request, issue ou nenhuma ação) mudou entre as quatro execuções em 4 das
20 unidades, mesmo entre execuções da mesma versão da biblioteca (Seção~\ref{subsec:rq6}).
Mudanças de decisão entre as reexecuções de outubro podem, portanto, refletir variação do modelo
tanto quanto mudanças no pipeline. Respostas em cache e pacotes de evidências tornam cada
decisão auditável, e as verificações e o build rejeitam formas inseguras conhecidas e propostas que não compilam,
mas a decisão em si não é reprodutível.

\emph{Julgamento do pesquisador.} Um único pesquisador executou o pipeline, escolheu o que
reexecutar e publicar e classificou os desfechos; as reexecuções favoreceram achados
próximos de um pull request. Algumas escolhas foram feitas depois de vistos os resultados
e são declaradas como desvios dos pré-registros: o refinamento do gabarito da comparação
e sua regra para correções já presentes na ponta do branch e, na validação
dinâmica, o rótulo para um patch que remove a chamada governada e o acréscimo
de 13 pull requests após os primeiros resultados. Os rótulos da validação dinâmica para um
caminho que o patch desliga, uma mensagem reportada que desapareceu, uma violação em
outro método do arquivo, um pull request só de testes e um relatório não mais
verificável também foram definidos depois de suas execuções; apenas pela regra registrada, 25 dos pull requests
contam como corrigidos e 13 como não corrigidos (33 e 4 com os rótulos posteriores;
Seção~\ref{subsec:dynval}). As respostas dos mantenedores são uma
verificação independente apenas para os casos publicados.

\emph{Comparação com agentes de programação.} Os agentes do primeiro bloco receberam apenas
o achado e o repositório, não a política do \system{} de que vale propor uma alteração
correta apenas de regra, e pararam no commit analisado
sem o porte para a ponta do branch; o segundo bloco recebe ambos. As 20 unidades, sorteadas com uma
semente registrada a partir dos pull requests e dos achados do estudo de caso deste trabalho,
representam este estudo, não o uso indevido de criptografia em Android em geral. Os
juízes são modelos de linguagem e nenhuma pessoa pontuou os cartões, pois
não havia revisores independentes disponíveis para ler os 491 cartões: o juiz
principal, Claude Opus 5.5, compartilha o provedor do braço Claude Code e desempata
divergências, e o GPT-6.1 Sol é o modelo de um dos braços Codex e compartilha o
provedor do modelo usado pelo \system{} e pelo outro braço Codex, de modo que um juiz pode
favorecer sua própria família. A referência de pontuação foi refinada duas vezes antes das
leituras reportadas, e a regra de pontuação da ponta do branch foi definida depois delas; apenas a
referência final é usada. Como uma issue ou nenhuma ação é aceitável em todas as
unidades, a medida recompensa a abstenção. O cegamento foi parcial; nas sondagens dos dois blocos, os
palpites do juiz principal sobre o braço foram palpites padrão, e não
identificações (Seção~\ref{subsec:rq6}). O prompt dos agentes, escrito a partir do pipeline, reproduz suas
etapas, mas não seus contextos separados, as verificações e novas tentativas que seu código
impõe, nem seu corpus de pesquisa; o Claude Code com Claude Sonnet 5.5 recusou
a versão que pedia um registro escrito de cada etapa, de modo que o SP2 usou um
prompt sem esse pedido. Os agentes rodaram em assinaturas com limite de 60 minutos
por execução. O sandbox dos agentes limitava as escritas ao seu checkout, mas não as
leituras: em duas execuções de um braço do primeiro bloco, uma busca por arquivos de
instruções listou caminhos de arquivos de outros repositórios da máquina; nenhuma
execução leu código de outro repositório nem de outra unidade da comparação.

\subsubsection{Validade de construto}

As observações iniciais são violações de regras, não vulnerabilidades exploráveis
confirmadas. O sucesso do build mede a compilação, não a corretude funcional,
a segurança ou a aceitação. Como nenhuma pessoa rotulou cada decisão
e cada patch, os resultados estabelecem viabilidade e modos de falha, não precisão
nem revocação. As regras podem ser incompletas ou incompatíveis com os provedores do Android, e
os monitores do RVSec foram traduzidos manualmente a partir do CrySL, de modo que uma regra errada pode moldar
os requisitos de evidência.

A validação dinâmica mede algo mais estreito que uma correção: se a
violação reportada deixou de aparecer no método modificado pelo patch. Um patch que apenas
tirasse a chamada do campo de visão do monitor passaria, e também um que
mantivesse o mesmo material de chave e mudasse apenas a forma como o monitor o vê,
como faz o FreeOTP \#524 (Seção~\ref{subsec:rq7}). Apenas pela regra registrada, um patch correto pode contar como não corrigido
quando outras restrições da especificação falham no mesmo método. Apenas os pull requests publicados foram reexecutados (as notas dos juízes
deixam o RVSec de fora), aplicando o diff ao commit sobre o qual o pipeline fez o build ou à
base do merge, não ao código mergeado. A validação depende da exploração: os novos
casos rodaram uma vez com Monkey e APE por 600 segundos (mais execuções apenas quando
inconclusivos ou não reproduzidos; os cinco casos anteriores, três vezes, o BitBanana,
quatro), e um caso conta como corrigido apenas quando o APK original mostra a violação
e o APK com o patch executa o método reportado sem ela, a menos que o patch tenha
removido a chamada governada ou desligado, por padrão, a condição que protege toda
chamada do método; outras mensagens da especificação que o APK original já mostrava
podem permanecer.
O método do FreeOTP, que no relatório original da campanha do grupo do
RVSec~\cite{costa2026effectiveness} apenas DroidBot, Humanoid e ARES alcançaram
e que Monkey e APE não alcançaram aqui, também foi reexecutado com o DroidBot, um dos
exploradores do relatório original (Humanoid e ARES não puderam rodar aqui;
Seção~\ref{subsec:dynval}).

\subsubsection{Validade externa}

O inventário é intencional: achados de Android analisados e contribuições
publicadas com artefatos rastreáveis, nem uma amostra aleatória de APKs nem
proporcional ao volume de alertas. Ele favorece achados com uma tentativa de
contribuição, pull requests abertos têm desfechos incompletos, e a política afeta a
aceitação. Os pull requests de nossas próprias filas do RVSec vieram de aplicações
cujos repositórios estavam ativos, haviam incorporado um pull request de um contribuidor
externo nos doze meses anteriores e não declaravam acordo de contribuidor nem regra de
contato prévio; isso favorece a incorporação em relação a uma amostra sem filtro. As contagens cobrem apenas achados do RVSec (não candidatos manuais, pesquisa
na web ou regras do Semgrep) e, com várias versões de APK, não constituem uma comparação
de detectores.
Apenas aplicações hospedadas no GitHub foram consideradas,
dado o caminho de publicação; as hospedadas em outras plataformas podem diferir. Os resultados podem não se generalizar para aplicações proprietárias, outras
linguagens ou outras famílias de API; três dos 49 pull requests tratam de
regras de gerenciamento de recursos ou de sincronização, e não de criptografia.
O \system{} aceita achados de outros detectores, mas nenhum está nessas contagens.

\subsubsection{Validade de conclusão}

As observações sustentam contagens descritivas e explicações de casos, não
afirmações inferenciais; nenhum teste estatístico é aplicado. Os 34 achados do
estudo de caso, os pull requests publicados, sua validação dinâmica e a
comparação têm denominadores diferentes e se sobrepõem parcialmente. A incorporação
não estabelece a corretude de segurança, nem o fechamento a incorretude; os resultados das
reexecuções se apoiam, em sua maioria, em uma única repetição, e as diferenças entre os braços
da comparação são reportadas como observadas, sem a afirmação de que se generalizam.

\subsection{Trabalhos Futuros}
\label{subsec:future}

Um estudo confirmatório deveria fixar suas hipóteses de antemão, sortear achados
de campanhas do RVSec e acrescentar duas revisões de segurança humanas independentes
por caso, de modo que a precisão das decisões de correção, a taxa de correções inseguras e
a adequação da abstenção possam ser medidas; a reexecução do detector deveria
usar várias repetições por APK e os exploradores do relatório original. Como
a exploração é o primeiro gargalo (Seção~\ref{sec:lessons}), uma exploração
guiada que priorize telas relacionadas às
especificações monitoradas poderia aumentar tanto os achados no código da aplicação quanto
os patches que a reexecução consegue confirmar.

As verificações do \system{} (validação semântica, build, porte, teto de custo) poderiam ser
oferecidas a agentes de programação como ferramentas ou como uma skill de agente, mantendo o julgamento separado
da correção e a publicação autorizada por uma pessoa; se isso preserva a qualidade das
decisões é uma questão para o protocolo da Seção~\ref{subsec:protocol}. O mesmo núcleo poderia rodar como uma GitHub
Action que analisa os achados do RVSec de cada build e anexa a alteração
proposta, ou o motivo de não haver nenhuma, onde os mantenedores já trabalham. Este estudo
não seguiu esse caminho: uma Action só roda em repositórios cujos mantenedores a instalam,
com suas próprias credenciais de modelo e um runner capaz de compilar a aplicação e hospedar
um emulador Android; como contribuidores externos, não podíamos executá-la nos projetos
estudados nem atribuir a eles o custo de suas execuções, e o estudo mediu, em vez disso, o que
os mantenedores fizeram com as contribuições submetidas a eles. Outros trabalhos incluem builds em contêiner e sem credenciais
para código de terceiros e, para TLS, um planejamento que considere todas as restrições
em \texttt{SSLContext.init}, já que a reexecução mostrou patches que removeram
o nome do protocolo enquanto o RVSec ainda reportava key managers, fontes de aleatoriedade ou
trust managers, com o pull request declarando quais restrições deixa
inalteradas.

\section{Conclusões}
\label{sec:conclusion}

Este artigo apresentou o \system{}, uma biblioteca Python que leva os achados do RVSec de
um relatório de execução sobre um APK a uma contribuição revisável: ela localiza a revisão
do código-fonte e o método, decide se a violação é real e se uma correção preserva o
comportamento da aplicação, constrói o patch, transporta-o para o branch atual e deixa
a publicação a cargo de uma pessoa. Sua premissa é que um relatório de uso indevido é uma
alegação a verificar, não um defeito a reparar, de modo que decidir não alterar o código
é um desfecho válido. O estudo do grupo do RVSec mede o que a verificação em tempo de
execução encontra~\cite{costa2026effectiveness}; este trabalho pergunta o que um
mantenedor deveria receber.

Os mantenedores incorporaram 21 dos 49 pull requests para achados do RVSec e fecharam 10
sem incorporá-los; 33 receberam resposta. Suas objeções raramente diziam respeito ao
código do patch: tratavam de se a regra importa naquela aplicação, da compatibilidade
com valores armazenados e de mudanças já feitas upstream, conhecimento que uma regra não
consegue codificar e que o \system{} agora lê para achados posteriores da mesma
especificação. Objeções ao texto ou à origem de um
pull request vieram apenas das submissões iniciais, com template
(Seção~\ref{subsec:form}).

A abstenção respondeu por boa parte do trabalho. No estudo de caso de outubro, a maioria
dos achados terminou sem correção gerada, 16 deles como propostas locais de issue e os
demais como falsos positivos, evidência insuficiente ou violação confirmada sem correção
segura (Tabela~\ref{tab:adjudication}). Nem toda mudança correta é desejada: seis pull
requests fechados substituíam um digest que o mantenedor via como checksum ou
identificador. Onde um patch foi publicado, a reexecução do RVSec no APK
com o patch deixou de mostrar a violação reportada no método alterado em 33 dos 38 pull
requests que ela pôde verificar (25 pela regra registrada), uma verificação limitada pelo
que os exploradores alcançam.

Nas 20 unidades da comparação, um agente de programação de uso geral não bastou sozinho. Recebendo apenas o achado e o
repositório, os agentes propuseram menos pull requests, com menos julgados aceitáveis;
recebendo o processo do \system{} como instruções, igualaram seus pull requests
aceitáveis na leitura completa, mas tiveram de três a cinco vezes mais julgados inaceitáveis.
Com o mesmo modelo e as mesmas instruções em um dos braços, a diferença é compatível com a
execução do processo em etapas separadas, verificadas por código, com a decisão de agir
tomada antes de qualquer edição. Uma execução custou cerca de US\$~0,10 no nível Flex
(US\$~0,19 a preço de tabela).

As evidências são descritivas e provêm de amostras pequenas e intencionais. Elas
sustentam o \system{} como um complemento ao RVSec orientado à revisão, e não a precisão
do detector nem a correção semântica e de segurança dos patches: uma violação que deixa
de aparecer mostra que o uso monitorado mudou, não que o comportamento foi preservado ou
que a aplicação se tornou mais segura. Os resultados também descrevem o \system{} operado
por seu autor, que escolheu o que reexecutar e aprovou cada submissão; eles não mostram
como ele se comportaria sem supervisão, o que seu projeto não pretende. Os próximos
passos decorrem das lições: uma exploração que alcance o código por trás de logins e
configurações, um estudo confirmatório com revisões de segurança independentes, e as
verificações do \system{} oferecidas como ferramentas a agentes de programação executados
em contêineres isolados (Seção~\ref{subsec:future}).

\section*{Disponibilidade dos Artefatos}

A biblioteca está disponível em \url{https://github.com/jpantonow/cognalys} e os
dados por trás de cada resultado em \url{https://github.com/jpantonow/cognalys-artefacts}. O repositório da
biblioteca contém o pacote, suas especificações e políticas de remediação, e seus
testes. O repositório de artefatos é organizado por questão de pesquisa: o retrato de
outubro de 2026, com as evidências, a adjudicação, o plano, os rastros de prompts, os
logs de build e as prévias de cada achado (\researchq{1}, \researchq{2},
\researchq{4}, \researchq{5}); os pull requests, as issues e as respostas dos
mantenedores coletados em 18 de novembro de 2026 (\researchq{3}); as unidades, as
execuções dos agentes, os cartões, o gabarito e as leituras dos juízes da comparação
(\researchq{6}); e os registros de build e os relatórios do RVSec da validação dinâmica
(\researchq{7}). Os pré-registros, com seus desvios declarados e os hashes dos prompts,
scripts e gabarito da comparação, estão incluídos com suas datas originais. As
credenciais foram excluídas.

\section*{Agradecimentos}

Agradecemos ao Prof.\ Rodrigo Bonifácio pela orientação deste trabalho, ao
Prof.\ Guilherme Novaes Ramos pela coorientação e a Pedro Henrique Teixeira Costa
pela campanha e pelas ferramentas do RVSec sobre as quais este trabalho se apoia.

\bibliographystyle{IEEEtran}
\bibliography{referencias}

\end{document}
